\documentclass{../../modelos/latex-shared/ifsc-ppc}

\begin{document}

% --- Capa ---
\capaPPC{Técnico em Administração Integrado}{2026}

% --- Sumário ---
\tableofcontents
\newpage

% =========================================================
\section{DADOS DA INSTITUIÇÃO}
Instituto Federal de Santa Catarina -- IFSC. \\
Instituído pela Lei nº 11.892, de 29 de dezembro de 2008. \\
Reitoria: Rua 14 de Julho, 150 -- Coqueiros -- Florianópolis -- Santa Catarina -- Brasil -- CEP 88.075-010 \\
Fone: +55 (48) 3877-9000 -- CNPJ: 11.402.887/0001-60.

% =========================================================
\section{DADOS DO CÂMPUS PROPONENTE}
\subsection{Câmpus:}
Garopaba

\subsection{Endereço e Telefone do Câmpus:}
Rua Maria Aparecida Barbosa, 153, Campo D'Una, Garopaba - CEP 88495-000 \\
Telefone: (48) 3254-7336

% =========================================================
\section{DADOS DOS RESPONSÁVEIS PELO PPC}
\subsection{Chefia DEPE / Departamento:}
Prof(a). [Aguardando confirmação da Chefia DEPE atual]

\subsection{Coordenador do curso/proponente:}
Profa. Fabiana de Agapito

\subsection{Equipe elaboradora do projeto de curso:}
Comissão de Reformulação do PPC Técnico em Administração (Portaria Câmpus Garopaba nº 2026/XXX)

\subsection{Aprovação no Câmpus:}
[Aguardando apreciação e aprovação pelo Colegiado de Câmpus -- Garopaba]

% =========================================================
\section{DADOS DO CURSO}

% --- TABELA 7: DADOS DO CURSO ---
\begin{table}[H]
\centering
\small
\begin{tabularx}{\linewidth}{|p{4.4cm}|X|p{3.8cm}|X|}
\hline
\rowcolor{cinzaTabela}
\multicolumn{4}{|l|}{\textcolor{ifscgreen}{\textbf{\large 7. Dados do Curso:}}} \\ \hline
\textcolor{ifscgreen}{\textbf{\textit{7.1. Nome do Curso:}}} & \multicolumn{3}{p{\dimexpr\linewidth-4.4cm-4\tabcolsep-3\arrayrulewidth\relax}|}{Curso Técnico em Administração.} \\ \hline
\textcolor{ifscgreen}{\textbf{\textit{7.2. Eixo Tecnológico:}}} & \multicolumn{3}{p{\dimexpr\linewidth-4.4cm-4\tabcolsep-3\arrayrulewidth\relax}|}{Gestão e negócios.} \\ \hline
\textcolor{ifscgreen}{\textbf{\textit{7.3 Forma de oferta}}} & \multicolumn{3}{p{\dimexpr\linewidth-4.4cm-4\tabcolsep-3\arrayrulewidth\relax}|}{Técnico Integrado ao Ensino Médio.\footnote{Conforme Lei nº 9.394, de 20 de dezembro de 1996, art. 36-C, alínea c e art. 46 do RDP.}} \\ \hline
\textcolor{ifscgreen}{\textbf{\textit{7.4. Modalidade:}}} & \multicolumn{3}{p{\dimexpr\linewidth-4.4cm-4\tabcolsep-3\arrayrulewidth\relax}|}{Presencial.} \\ \hline
\textcolor{ifscgreen}{\textbf{\textit{7.5. Diplomação do Egresso:}}} & \multicolumn{3}{p{\dimexpr\linewidth-4.4cm-4\tabcolsep-3\arrayrulewidth\relax}|}{Técnico em Administração;} \\ \hline
\textcolor{ifscgreen}{\textbf{\textit{7.6. CH Total:}}} & \multicolumn{3}{p{\dimexpr\linewidth-4.4cm-4\tabcolsep-3\arrayrulewidth\relax}|}{3260 horas} \\ \hline
\textcolor{ifscgreen}{\textbf{7.6.1 CH Aulas presenciais}} & 3260 horas & \textcolor{ifscgreen}{\textbf{7.6.2. CH Aulas EaD:}} & 0 h \\ \hline
\textcolor{ifscgreen}{\textbf{7.6.3 CH Estágio}} & \multicolumn{3}{p{\dimexpr\linewidth-4.4cm-4\tabcolsep-3\arrayrulewidth\relax}|}{0 horas (não obrigatório).} \\ \hline
\textcolor{ifscgreen}{\textbf{7.6.4. Demais CH previstas}} & \multicolumn{3}{p{\dimexpr\linewidth-4.4cm-4\tabcolsep-3\arrayrulewidth\relax}|}{Não se aplica.} \\ \hline
\end{tabularx}
\caption{Dados Gerais do Curso Técnico em Administração}
\end{table}

\vspace{0.3cm}

% --- TABELA 8: DADOS DA OFERTA ---
\begin{table}[H]
\centering
\small
\begin{tabularx}{\linewidth}{|p{4.4cm}|X|p{3.8cm}|X|}
\hline
\rowcolor{cinzaTabela}
\multicolumn{4}{|l|}{\textcolor{ifscgreen}{\textbf{\large 8. Dados da Oferta:}}} \\ \hline
\textcolor{ifscgreen}{\textbf{\textit{8.1. Local da Oferta:}}} & \multicolumn{3}{p{\dimexpr\linewidth-4.4cm-4\tabcolsep-3\arrayrulewidth\relax}|}{IFSC Câmpus Garopaba. Rua Maria Aparecida Barbosa, 153 -- Campo D'Una, Garopaba - SC, CEP 88495-000, (48) 3254-7336.} \\ \hline
\textcolor{ifscgreen}{\textbf{\textit{8.2. Vagas por Turma:}}} & \multicolumn{3}{p{\dimexpr\linewidth-4.4cm-4\tabcolsep-3\arrayrulewidth\relax}|}{40 vagas por turma.} \\ \hline
\textcolor{ifscgreen}{\textbf{\textit{8.3. Vagas totais Anuais:}}} & \multicolumn{3}{p{\dimexpr\linewidth-4.4cm-4\tabcolsep-3\arrayrulewidth\relax}|}{40 vagas anuais.} \\ \hline
\textcolor{ifscgreen}{\textbf{\textit{8.4. Periodicidade de ingresso no curso}}} & \multicolumn{3}{p{\dimexpr\linewidth-4.4cm-4\tabcolsep-3\arrayrulewidth\relax}|}{Anual.} \\ \hline
\textcolor{ifscgreen}{\textbf{\textit{8.5. Turno:}}} & \multicolumn{3}{p{\dimexpr\linewidth-4.4cm-4\tabcolsep-3\arrayrulewidth\relax}|}{Vespertino -- atividades no contraturno até duas vezes por semana.} \\ \hline
\textcolor{ifscgreen}{\textbf{\textit{8.6. Integralização (Mínimo):}}} & 06 semestres (3 anos) & \textcolor{ifscgreen}{\textbf{\textit{8.7 Integralização (Máximo):}}} & 10 semestres (5 anos) \\ \hline
\textcolor{ifscgreen}{\textbf{\textit{8.8. Regime de matrícula:}}} & \multicolumn{3}{p{\dimexpr\linewidth-4.4cm-4\tabcolsep-3\arrayrulewidth\relax}|}{Matrícula por ano letivo / bloco de unidades curriculares.} \\ \hline
\textcolor{ifscgreen}{\textbf{\textit{8.9. Forma de Ingresso:}}} & \multicolumn{3}{p{\dimexpr\linewidth-4.4cm-4\tabcolsep-3\arrayrulewidth\relax}|}{Exame de Seleção / Processo Seletivo IFSC.} \\ \hline
\textcolor{ifscgreen}{\textbf{\textit{8.10. Previsão de início da Oferta}}} & \multicolumn{3}{p{\dimexpr\linewidth-4.4cm-4\tabcolsep-3\arrayrulewidth\relax}|}{2026/1} \\ \hline
\end{tabularx}
\caption{Dados da Oferta do Curso Técnico em Administração}
\end{table}

\subsection{Requisito de Acesso:}
Conclusão do Ensino Fundamental (ou equivalente legalmente reconhecido).

\subsection{Legislação Aplicada ao Curso:}
\begin{itemize}
    \item Lei nº 9.394/1996 (Diretrizes e Bases da Educação Nacional -- LDB);
    \item Lei nº 11.892/2008 (Criação do Instituto Federal de Santa Catarina);
    \item Resolução CNE/CP nº 1/2021 (Diretrizes Curriculares Nacionais para a Educação Profissional e Tecnológica);
    \item Catálogo Nacional de Cursos Técnicos (CNCT/MEC) -- 4ª Edição;
    \item Resolução CONSUP/IFSC nº 142/2025 (Diretrizes Curriculares do Ensino Médio Integrado no IFSC);
    \item Regulamento Didático-Pedagógico do IFSC (RDP).
\end{itemize}

\subsection{Justificativa da Oferta do Curso no Câmpus Garopaba:}
A região abrangida pelo Câmpus Garopaba apresenta vocação econômica fortemente vinculada ao comércio, turismo, prestação de serviços e desenvolvimento sustentável local. A oferta do Curso Técnico em Administração Integrado ao Ensino Médio visa capacitar jovens da comunidade para atuar na gestão eficiente, empreendedorismo socioambiental, logística, finanças e processos organizacionais de micro, pequenas e médias empresas regionais.

\subsection{Público-Alvo:}
Estudantes egressos do Ensino Fundamental que buscam concluir a formação do Ensino Médio aliada a uma qualificação profissional técnica no eixo de Gestão e Negócios.

\subsection{Objetivos do Curso:}
\subsubsection{Objetivo Geral:}
Formar técnicos em Administração aptos a executar, analisar e propor melhorias em processos organizacionais, administrativos, financeiros, mercadológicos e de pessoas, comprometidos com a ética, sustentabilidade socioambiental e inovação.

\subsubsection{Objetivos Específicos:}
\begin{itemize}
    \item Desenvolver competências técnicas e científicas para operacionalização de rotinas administrativas, contábeis e financeiras;
    \item Capacitar para a aplicação de ferramentas de gestão da qualidade, marketing, empreendedorismo e logística;
    \item Promover a visão crítica sobre o papel das organizações no desenvolvimento sustentável regional;
    \item Integrar os conhecimentos das disciplinas da Formação Geral ao contexto prático profissional do Núcleo Politécnico e da Formação Técnica.
\end{itemize}

\subsection{Perfil Profissional do Egresso:}
O Técnico em Administração egresso do IFSC Câmpus Garopaba executa operações administrativas relativas aos processos de gestão de pessoas, materiais e patrimônio, compras, vendas, finanças, custos e marketing. Domina ferramentas de informática aplicada à gestão, elabora relatórios e atua no suporte aos processos decisórios organizacionais com responsabilidade social e ambiental.

\subsection{Áreas/Campos de Atuação do Egresso:}
Empresas públicas e privadas do comércio, indústria e serviços; organizações não governamentais (ONGs); cooperativas; startups e empreendimentos próprios; órgãos da administração pública direta e indireta.

% =========================================================
\section{ESTRUTURA CURRICULAR DO CURSO}
\subsection{Metodologia de Desenvolvimento Pedagógico:}
O curso orienta-se pelos princípios da integração curricular, articulando ciência, cultura, trabalho e tecnologia. As atividades pedagógicas incluem aulas expositivo-dialogadas, estudos de caso, atividades em laboratório de informática aplicada, pesquisas de campo, oficinas de integração e projetos integradores.

\subsection{Síntese da Matriz Curricular (Distribuição por Blocos):}
\begin{table}[H]
\centering
\begin{tabular}{|l|c|c|}
\hline
\textbf{Bloco Curricular (Res. 142/2025)} & \textbf{Carga Horária Mínima} & \textbf{Percentual Total} \\ \hline
Formação Geral (BNCC + Diversificada) & 2.140 h & 65,6\% \\ \hline
Formação Técnica Profissional (Administração) & 1.000 h & 30,7\% \\ \hline
Núcleo Politécnico Comum (Oficinas I, II e III) & 120 h & 3,7\% \\ \hline
\textbf{Carga Horária Total do Curso} & \textbf{3.260 h} & \textbf{100,0\%} \\ \hline
\end{tabular}
\caption{Matriz Curricular Resumida -- Técnico em Administração 2026}
\end{table}

% =========================================================
% EMENTÁRIO COMPLETO DO CURSO TÉCNICO EM ADMINISTRAÇÃO (PPC 2026)
% Formatado no Padrão Oficial IFSC (Resolução CONSUP/IFSC nº 142/2025 e CEPE nº 28/2025)
% =========================================================

\subsection{Matriz Curricular Detalhada por Anos e Blocos}

% --- TABELA 1º ANO ---
\begin{table}[H]
\centering
\small
\begin{tabularx}{\linewidth}{|c|X|c|c|}
\hline
\rowcolor{ifscgreen!20}
\multicolumn{4}{|c|}{\textbf{\large 1º ANO LETIVO (SEMESTRES 1 E 2)}} \\ \hline
\rowcolor{cinzaTabela}
\textbf{Ano} & \textbf{Unidade Curricular (UC)} & \textbf{Bloco Curricular} & \textbf{CH (h)} \\ \hline
1º Ano & Língua Portuguesa e Literatura I & Formação Geral & 80 h \\ \hline
1º Ano & Matemática I & Formação Geral & 80 h \\ \hline
1º Ano & História I & Formação Geral & 80 h \\ \hline
1º Ano & Geografia I & Formação Geral & 80 h \\ \hline
1º Ano & Biologia I & Formação Geral & 80 h \\ \hline
1º Ano & Física I & Formação Geral & 80 h \\ \hline
1º Ano & Química I & Formação Geral & 80 h \\ \hline
1º Ano & Filosofia I & Formação Geral & 80 h \\ \hline
1º Ano & Educação Física I & Formação Geral & 80 h \\ \hline
1º Ano & Espanhol & Formação Geral & 80 h \\ \hline
1º Ano & Introdução à Administração & Formação Técnica & 80 h \\ \hline
1º Ano & Informática Aplicada & Formação Técnica & 60 h \\ \hline
1º Ano & Organização e Processos & Formação Técnica & 80 h \\ \hline
1º Ano & Sociedade e Trabalho & Formação Técnica & 60 h \\ \hline
1º Ano & Oficina de Integração I & Núcleo Politécnico & 40 h \\ \hline
\multicolumn{3}{|r|}{\textbf{Subtotal 1º Ano (720h FG + 280h FT + 40h NP):}} & \textbf{1.040 h} \\ \hline
\end{tabularx}
\caption{Matriz Curricular -- 1º Ano Letivo}
\end{table}

% --- TABELA 2º ANO ---
\begin{table}[H]
\centering
\small
\begin{tabularx}{\linewidth}{|c|X|c|c|}
\hline
\rowcolor{ifscgreen!20}
\multicolumn{4}{|c|}{\textbf{\large 2º ANO LETIVO (SEMESTRES 1 E 2)}} \\ \hline
\rowcolor{cinzaTabela}
\textbf{Ano} & \textbf{Unidade Curricular (UC)} & \textbf{Bloco Curricular} & \textbf{CH (h)} \\ \hline
2º Ano & Língua Portuguesa e Literatura II & Formação Geral & 80 h \\ \hline
2º Ano & Matemática II & Formação Geral & 80 h \\ \hline
2º Ano & História II & Formação Geral & 40 h \\ \hline
2º Ano & Geografia II & Formação Geral & 80 h \\ \hline
2º Ano & Biologia II & Formação Geral & 80 h \\ \hline
2º Ano & Física II & Formação Geral & 40 h \\ \hline
2º Ano & Química II & Formação Geral & 40 h \\ \hline
2º Ano & Sociologia I & Formação Geral & 80 h \\ \hline
2º Ano & Artes & Formação Geral & 80 h \\ \hline
2º Ano & Educação Física II & Formação Geral & 80 h \\ \hline
2º Ano & Matemática para Administração & Formação Técnica & 80 h \\ \hline
2º Ano & Gestão de Pessoas e Relações no Trabalho & Formação Técnica & 80 h \\ \hline
2º Ano & Gestão de Marketing I & Formação Técnica & 80 h \\ \hline
2º Ano & Gestão Financeira & Formação Técnica & 80 h \\ \hline
2º Ano & Gestão de Marketing II & Formação Técnica & 80 h \\ \hline
2º Ano & Empreendedorismo I & Formação Técnica & 80 h \\ \hline
2º Ano & Oficina de Integração II & Núcleo Politécnico & 40 h \\ \hline
\multicolumn{3}{|r|}{\textbf{Subtotal 2º Ano (660h FG + 480h FT + 40h NP):}} & \textbf{1.180 h} \\ \hline
\end{tabularx}
\caption{Matriz Curricular -- 2º Ano Letivo}
\end{table}

\newpage

% --- TABELA 3º ANO ---
\begin{table}[H]
\centering
\small
\begin{tabularx}{\linewidth}{|c|X|c|c|}
\hline
\rowcolor{ifscgreen!20}
\multicolumn{4}{|c|}{\textbf{\large 3º ANO LETIVO (SEMESTRES 1 E 2)}} \\ \hline
\rowcolor{cinzaTabela}
\textbf{Ano} & \textbf{Unidade Curricular (UC)} & \textbf{Bloco Curricular} & \textbf{CH (h)} \\ \hline
3º Ano & Língua Portuguesa e Literatura III & Formação Geral & 120 h \\ \hline
3º Ano & Matemática III & Formação Geral & 80 h \\ \hline
3º Ano & História III & Formação Geral & 40 h \\ \hline
3º Ano & Geografia III & Formação Geral & 40 h \\ \hline
3º Ano & Filosofia II & Formação Geral & 40 h \\ \hline
3º Ano & Sociologia II & Formação Geral & 40 h \\ \hline
3º Ano & Física III & Formação Geral & 80 h \\ \hline
3º Ano & Química III & Formação Geral & 80 h \\ \hline
3º Ano & Inglês & Formação Geral & 80 h \\ \hline
3º Ano & Artes II & Formação Geral & 60 h \\ \hline
3º Ano & Gestão de Operações e Qualidade & Formação Técnica & 80 h \\ \hline
3º Ano & Empreendedorismo II (Plano de Negócios) & Formação Técnica & 80 h \\ \hline
3º Ano & Responsabilidade Socioambiental e Sustentabilidade & Formação Técnica & 80 h \\ \hline
3º Ano & Oficina de Integração III (Projeto Integrador) & Núcleo Politécnico & 40 h \\ \hline
\multicolumn{3}{|r|}{\textbf{Subtotal 3º Ano (760h FG + 240h FT + 40h NP):}} & \textbf{1.040 h} \\ \hline
\end{tabularx}
\caption{Matriz Curricular -- 3º Ano Letivo}
\end{table}

% --- QUADRO SÍNTESE FINAL ---
\begin{table}[H]
\centering
\small
\begin{tabularx}{\linewidth}{|X|c|c|}
\hline
\rowcolor{ifscgreen!20}
\multicolumn{3}{|c|}{\textbf{\large QUADRO SÍNTESE DA CARGA HORÁRIA TOTAL DO CURSO}} \\ \hline
\rowcolor{cinzaTabela}
\textbf{Bloco Curricular} & \textbf{Carga Horária (horas-relógio)} & \textbf{Percentual Total} \\ \hline
Formação Geral (Base Nacional Comum Curricular) & 2.140 h & 65,6 \% \\ \hline
Formação Técnica Profissional (Eixo Gestão e Negócios) & 1.000 h & 30,7 \% \\ \hline
Núcleo Politécnico Comum (Oficinas de Integração I, II e III) & 120 h & 3,7 \% \\ \hline
\hline
\rowcolor{ifscgreen!10}
\textbf{CARGA HORÁRIA TOTAL GERAL DO CURSO} & \textbf{3.260 h} & \textbf{100,0 \%} \\ \hline
\end{tabularx}
\caption{Síntese da Carga Horária do Técnico em Administração (PPC 2026)}
\end{table}

\newpage
\subsection{Ementário das Unidades Curriculares (Formato Oficial em Tabela)}

% ---------------------------------------------------------
% TABELA EMENTA 1: LÍNGUA PORTUGUESA E LITERATURA I
% ---------------------------------------------------------
\begin{table}[H]
\centering
\small
\begin{tabularx}{\linewidth}{|p{3.4cm}|X|p{3.5cm}|X|}
\hline
\multicolumn{4}{|>{\centering\arraybackslash\cellcolor{ifscgreen!20}}p{\dimexpr\linewidth-2\tabcolsep-2\arrayrulewidth\relax}|}{\textbf{\large Unidade Curricular: LÍNGUA PORTUGUESA E LITERATURA I}} \\ \hline
\textbf{CH Total*:} & 80 horas & \textbf{Ano / Semestre*:} & 1º Ano (Semestres 1 e 2) \\ \hline
\textbf{CH* Prática:} & 20 horas & \textbf{CH EaD*:} & não se aplica \\ \hline
\multicolumn{2}{|l|}{\textbf{CH com Divisão de Turma*:}} & \multicolumn{2}{l|}{não se aplica} \\ \hline
\multicolumn{4}{|p{\dimexpr\linewidth-2\tabcolsep-2\arrayrulewidth\relax}|}{\cellcolor{cinzaTabela}\textbf{Objetivos:}} \\ \hline
\multicolumn{4}{|p{\dimexpr\linewidth-2\tabcolsep-2\arrayrulewidth\relax}|}{
\begin{itemize}[leftmargin=*,noitemsep,topsep=2pt]
    \item Compreender os aspectos fundamentais da linguagem verbal e não verbal no contexto acadêmico e corporativo;
    \item Desenvolver a leitura crítica e a produção textual de gêneros discursivos diversos;
\end{itemize}
} \\ \hline
\multicolumn{4}{|p{\dimexpr\linewidth-2\tabcolsep-2\arrayrulewidth\relax}|}{\cellcolor{cinzaTabela}\textbf{Conteúdos:}} \\ \hline
\multicolumn{4}{|p{\dimexpr\linewidth-2\tabcolsep-2\arrayrulewidth\relax}|}{
\begin{itemize}[leftmargin=*,noitemsep,topsep=2pt]
    \item Leitura, interpretação e produção de textos acadêmicos e técnicos.
    \item Gramática aplicada e variação linguística.
    \item Introdução aos estudos literários e produções culturais.
\end{itemize}
} \\ \hline
\multicolumn{4}{|p{\dimexpr\linewidth-2\tabcolsep-2\arrayrulewidth\relax}|}{\cellcolor{cinzaTabela}\textbf{Metodologia de Abordagem:}} \\ \hline
\multicolumn{4}{|p{\dimexpr\linewidth-2\tabcolsep-2\arrayrulewidth\relax}|}{Aulas expositivo-dialogadas, oficinas de leitura e produção textual com suporte de recursos multimídia. Avaliação processual mediante produções escritas e participação ativa.} \\ \hline
\multicolumn{4}{|p{\dimexpr\linewidth-2\tabcolsep-2\arrayrulewidth\relax}|}{\cellcolor{cinzaTabela}\textbf{Bibliografia Básica:}} \\ \hline
\multicolumn{4}{|p{\dimexpr\linewidth-2\tabcolsep-2\arrayrulewidth\relax}|}{1. CEREJA, William Roberto; MAGALHÃES, Cleyvis. Gramática Refletiva. 4. ed. São Paulo: Atual, 2020.
2. FIORIN, José Luiz; SAVIOLI, Francisco Platão. Para Entender o Texto. 17. ed. São Paulo: Ática, 2018.} \\ \hline
\multicolumn{4}{|p{\dimexpr\linewidth-2\tabcolsep-2\arrayrulewidth\relax}|}{\cellcolor{cinzaTabela}\textbf{Bibliografia Complementar:}} \\ \hline
\multicolumn{4}{|p{\dimexpr\linewidth-2\tabcolsep-2\arrayrulewidth\relax}|}{1. KOCH, Ingedore Villaça. A Coesão Textual. 22. ed. São Paulo: Contexto, 2019.} \\ \hline
\end{tabularx}
\end{table}

\vspace{0.3cm}


% ---------------------------------------------------------
% TABELA EMENTA 2: MATEMÁTICA I
% ---------------------------------------------------------
\begin{table}[H]
\centering
\small
\begin{tabularx}{\linewidth}{|p{3.4cm}|X|p{3.5cm}|X|}
\hline
\multicolumn{4}{|>{\centering\arraybackslash\cellcolor{ifscgreen!20}}p{\dimexpr\linewidth-2\tabcolsep-2\arrayrulewidth\relax}|}{\textbf{\large Unidade Curricular: MATEMÁTICA I}} \\ \hline
\textbf{CH Total*:} & 80 horas & \textbf{Ano / Semestre*:} & 1º Ano (Semestres 1 e 2) \\ \hline
\textbf{CH* Prática:} & 20 horas & \textbf{CH EaD*:} & não se aplica \\ \hline
\multicolumn{2}{|l|}{\textbf{CH com Divisão de Turma*:}} & \multicolumn{2}{l|}{não se aplica} \\ \hline
\multicolumn{4}{|p{\dimexpr\linewidth-2\tabcolsep-2\arrayrulewidth\relax}|}{\cellcolor{cinzaTabela}\textbf{Objetivos:}} \\ \hline
\multicolumn{4}{|p{\dimexpr\linewidth-2\tabcolsep-2\arrayrulewidth\relax}|}{
\begin{itemize}[leftmargin=*,noitemsep,topsep=2pt]
    \item Compreender e aplicar conceitos de conjuntos, funções e razões nas resoluções de problemas cotidianos e técnicos.
\end{itemize}
} \\ \hline
\multicolumn{4}{|p{\dimexpr\linewidth-2\tabcolsep-2\arrayrulewidth\relax}|}{\cellcolor{cinzaTabela}\textbf{Conteúdos:}} \\ \hline
\multicolumn{4}{|p{\dimexpr\linewidth-2\tabcolsep-2\arrayrulewidth\relax}|}{
\begin{itemize}[leftmargin=*,noitemsep,topsep=2pt]
    \item Teoria dos Conjuntos e Conjuntos Numéricos.
    \item Funções do 1º e 2º Grau e Aplicações.
    \item Razão, Proporção, Regra de Três e Porcentagem.
\end{itemize}
} \\ \hline
\multicolumn{4}{|p{\dimexpr\linewidth-2\tabcolsep-2\arrayrulewidth\relax}|}{\cellcolor{cinzaTabela}\textbf{Metodologia de Abordagem:}} \\ \hline
\multicolumn{4}{|p{\dimexpr\linewidth-2\tabcolsep-2\arrayrulewidth\relax}|}{Aulas expositivas e resolução de listas de exercícios orientados. Avaliações escritas individuais e trabalhos práticos.} \\ \hline
\multicolumn{4}{|p{\dimexpr\linewidth-2\tabcolsep-2\arrayrulewidth\relax}|}{\cellcolor{cinzaTabela}\textbf{Bibliografia Básica:}} \\ \hline
\multicolumn{4}{|p{\dimexpr\linewidth-2\tabcolsep-2\arrayrulewidth\relax}|}{1. DANTE, Luiz Roberto. Matemática: Contexto e Aplicações. 3. ed. São Paulo: Ática, 2020.} \\ \hline
\multicolumn{4}{|p{\dimexpr\linewidth-2\tabcolsep-2\arrayrulewidth\relax}|}{\cellcolor{cinzaTabela}\textbf{Bibliografia Complementar:}} \\ \hline
\multicolumn{4}{|p{\dimexpr\linewidth-2\tabcolsep-2\arrayrulewidth\relax}|}{1. IEZZI, Gelson et al. Fundamentos de Matemática Elementar. 9. ed. São Paulo: Atual, 2018.} \\ \hline
\end{tabularx}
\end{table}

\vspace{0.3cm}


% ---------------------------------------------------------
% TABELA EMENTA 3: HISTÓRIA I
% ---------------------------------------------------------
\begin{table}[H]
\centering
\small
\begin{tabularx}{\linewidth}{|p{3.4cm}|X|p{3.5cm}|X|}
\hline
\multicolumn{4}{|>{\centering\arraybackslash\cellcolor{ifscgreen!20}}p{\dimexpr\linewidth-2\tabcolsep-2\arrayrulewidth\relax}|}{\textbf{\large Unidade Curricular: HISTÓRIA I}} \\ \hline
\textbf{CH Total*:} & 80 horas & \textbf{Ano / Semestre*:} & 1º Ano (Semestres 1 e 2) \\ \hline
\textbf{CH* Prática:} & 0 horas & \textbf{CH EaD*:} & não se aplica \\ \hline
\multicolumn{2}{|l|}{\textbf{CH com Divisão de Turma*:}} & \multicolumn{2}{l|}{não se aplica} \\ \hline
\multicolumn{4}{|p{\dimexpr\linewidth-2\tabcolsep-2\arrayrulewidth\relax}|}{\cellcolor{cinzaTabela}\textbf{Objetivos:}} \\ \hline
\multicolumn{4}{|p{\dimexpr\linewidth-2\tabcolsep-2\arrayrulewidth\relax}|}{
\begin{itemize}[leftmargin=*,noitemsep,topsep=2pt]
    \item Analisar as transformações históricas da Antiguidade ao Período Colonial e suas implicações socioeconômicas.
\end{itemize}
} \\ \hline
\multicolumn{4}{|p{\dimexpr\linewidth-2\tabcolsep-2\arrayrulewidth\relax}|}{\cellcolor{cinzaTabela}\textbf{Conteúdos:}} \\ \hline
\multicolumn{4}{|p{\dimexpr\linewidth-2\tabcolsep-2\arrayrulewidth\relax}|}{
\begin{itemize}[leftmargin=*,noitemsep,topsep=2pt]
    \item Antiguidade Clássica e Feudalismo.
    \item Mercantilismo e Expansão Marítima.
    \item Brasil Colonial e Formação Social.
\end{itemize}
} \\ \hline
\multicolumn{4}{|p{\dimexpr\linewidth-2\tabcolsep-2\arrayrulewidth\relax}|}{\cellcolor{cinzaTabela}\textbf{Metodologia de Abordagem:}} \\ \hline
\multicolumn{4}{|p{\dimexpr\linewidth-2\tabcolsep-2\arrayrulewidth\relax}|}{Aulas dialogadas com apoio de documentos históricos e filmes. Avaliações escritas e seminários.} \\ \hline
\multicolumn{4}{|p{\dimexpr\linewidth-2\tabcolsep-2\arrayrulewidth\relax}|}{\cellcolor{cinzaTabela}\textbf{Bibliografia Básica:}} \\ \hline
\multicolumn{4}{|p{\dimexpr\linewidth-2\tabcolsep-2\arrayrulewidth\relax}|}{1. VAINFAS, Ronaldo et al. História Conectada. 2. ed. São Paulo: Saraiva, 2020.} \\ \hline
\multicolumn{4}{|p{\dimexpr\linewidth-2\tabcolsep-2\arrayrulewidth\relax}|}{\cellcolor{cinzaTabela}\textbf{Bibliografia Complementar:}} \\ \hline
\multicolumn{4}{|p{\dimexpr\linewidth-2\tabcolsep-2\arrayrulewidth\relax}|}{1. HOLANDA, Sérgio Buarque de. Raízes do Brasil. São Paulo: Companhia das Letras, 2016.} \\ \hline
\end{tabularx}
\end{table}

\vspace{0.3cm}


% ---------------------------------------------------------
% TABELA EMENTA 4: GEOGRAFIA I
% ---------------------------------------------------------
\begin{table}[H]
\centering
\small
\begin{tabularx}{\linewidth}{|p{3.4cm}|X|p{3.5cm}|X|}
\hline
\multicolumn{4}{|>{\centering\arraybackslash\cellcolor{ifscgreen!20}}p{\dimexpr\linewidth-2\tabcolsep-2\arrayrulewidth\relax}|}{\textbf{\large Unidade Curricular: GEOGRAFIA I}} \\ \hline
\textbf{CH Total*:} & 80 horas & \textbf{Ano / Semestre*:} & 1º Ano (Semestres 1 e 2) \\ \hline
\textbf{CH* Prática:} & 20 horas & \textbf{CH EaD*:} & não se aplica \\ \hline
\multicolumn{2}{|l|}{\textbf{CH com Divisão de Turma*:}} & \multicolumn{2}{l|}{não se aplica} \\ \hline
\multicolumn{4}{|p{\dimexpr\linewidth-2\tabcolsep-2\arrayrulewidth\relax}|}{\cellcolor{cinzaTabela}\textbf{Objetivos:}} \\ \hline
\multicolumn{4}{|p{\dimexpr\linewidth-2\tabcolsep-2\arrayrulewidth\relax}|}{
\begin{itemize}[leftmargin=*,noitemsep,topsep=2pt]
    \item Compreender as dinâmicas do espaço geográfico, cartografia e geopolítica ambiental.
\end{itemize}
} \\ \hline
\multicolumn{4}{|p{\dimexpr\linewidth-2\tabcolsep-2\arrayrulewidth\relax}|}{\cellcolor{cinzaTabela}\textbf{Conteúdos:}} \\ \hline
\multicolumn{4}{|p{\dimexpr\linewidth-2\tabcolsep-2\arrayrulewidth\relax}|}{
\begin{itemize}[leftmargin=*,noitemsep,topsep=2pt]
    \item Cartografia e Orientação no Espaço.
    \item Geomorfologia e Dinâmicas Naturais.
    \item Geopolítica e Meio Ambiente.
\end{itemize}
} \\ \hline
\multicolumn{4}{|p{\dimexpr\linewidth-2\tabcolsep-2\arrayrulewidth\relax}|}{\cellcolor{cinzaTabela}\textbf{Metodologia de Abordagem:}} \\ \hline
\multicolumn{4}{|p{\dimexpr\linewidth-2\tabcolsep-2\arrayrulewidth\relax}|}{Aulas expositivas, análise de mapas e saídas de campo. Avaliações escritas e trabalhos práticos de cartografia.} \\ \hline
\multicolumn{4}{|p{\dimexpr\linewidth-2\tabcolsep-2\arrayrulewidth\relax}|}{\cellcolor{cinzaTabela}\textbf{Bibliografia Básica:}} \\ \hline
\multicolumn{4}{|p{\dimexpr\linewidth-2\tabcolsep-2\arrayrulewidth\relax}|}{1. LUCCI, Elian Alabi et al. Território e Sociedade no Mundo Globalizado. São Paulo: Saraiva, 2020.} \\ \hline
\multicolumn{4}{|p{\dimexpr\linewidth-2\tabcolsep-2\arrayrulewidth\relax}|}{\cellcolor{cinzaTabela}\textbf{Bibliografia Complementar:}} \\ \hline
\multicolumn{4}{|p{\dimexpr\linewidth-2\tabcolsep-2\arrayrulewidth\relax}|}{1. SANTOS, Milton. A Natureza do Espaço. São Paulo: EDUSP, 2014.} \\ \hline
\end{tabularx}
\end{table}

\vspace{0.3cm}


% ---------------------------------------------------------
% TABELA EMENTA 5: BIOLOGIA I
% ---------------------------------------------------------
\begin{table}[H]
\centering
\small
\begin{tabularx}{\linewidth}{|p{3.4cm}|X|p{3.5cm}|X|}
\hline
\multicolumn{4}{|>{\centering\arraybackslash\cellcolor{ifscgreen!20}}p{\dimexpr\linewidth-2\tabcolsep-2\arrayrulewidth\relax}|}{\textbf{\large Unidade Curricular: BIOLOGIA I}} \\ \hline
\textbf{CH Total*:} & 80 horas & \textbf{Ano / Semestre*:} & 1º Ano (Semestres 1 e 2) \\ \hline
\textbf{CH* Prática:} & 20 horas & \textbf{CH EaD*:} & não se aplica \\ \hline
\multicolumn{2}{|l|}{\textbf{CH com Divisão de Turma*:}} & \multicolumn{2}{l|}{não se aplica} \\ \hline
\multicolumn{4}{|p{\dimexpr\linewidth-2\tabcolsep-2\arrayrulewidth\relax}|}{\cellcolor{cinzaTabela}\textbf{Objetivos:}} \\ \hline
\multicolumn{4}{|p{\dimexpr\linewidth-2\tabcolsep-2\arrayrulewidth\relax}|}{
\begin{itemize}[leftmargin=*,noitemsep,topsep=2pt]
    \item Compreender a estrutura celular, processos bioquímicos e ecológicos.
\end{itemize}
} \\ \hline
\multicolumn{4}{|p{\dimexpr\linewidth-2\tabcolsep-2\arrayrulewidth\relax}|}{\cellcolor{cinzaTabela}\textbf{Conteúdos:}} \\ \hline
\multicolumn{4}{|p{\dimexpr\linewidth-2\tabcolsep-2\arrayrulewidth\relax}|}{
\begin{itemize}[leftmargin=*,noitemsep,topsep=2pt]
    \item Citologia e Bioquímica Celular.
    \item Ecologia e Ecossistemas Regionais.
    \item Biodiversidade e Sustentabilidade.
\end{itemize}
} \\ \hline
\multicolumn{4}{|p{\dimexpr\linewidth-2\tabcolsep-2\arrayrulewidth\relax}|}{\cellcolor{cinzaTabela}\textbf{Metodologia de Abordagem:}} \\ \hline
\multicolumn{4}{|p{\dimexpr\linewidth-2\tabcolsep-2\arrayrulewidth\relax}|}{Aulas expositivas e práticas de laboratório. Avaliações escritas e relatórios de aula prática.} \\ \hline
\multicolumn{4}{|p{\dimexpr\linewidth-2\tabcolsep-2\arrayrulewidth\relax}|}{\cellcolor{cinzaTabela}\textbf{Bibliografia Básica:}} \\ \hline
\multicolumn{4}{|p{\dimexpr\linewidth-2\tabcolsep-2\arrayrulewidth\relax}|}{1. AMABIS, José Mariano; MARTHO, Gilberto Rodrigues. Biologia. 4. ed. São Paulo: Moderna, 2020.} \\ \hline
\multicolumn{4}{|p{\dimexpr\linewidth-2\tabcolsep-2\arrayrulewidth\relax}|}{\cellcolor{cinzaTabela}\textbf{Bibliografia Complementar:}} \\ \hline
\multicolumn{4}{|p{\dimexpr\linewidth-2\tabcolsep-2\arrayrulewidth\relax}|}{1. CAMPBELL, Neil A. et al. Biologia de Campbell. 10. ed. Porto Alegre: Artmed, 2015.} \\ \hline
\end{tabularx}
\end{table}

\vspace{0.3cm}


% ---------------------------------------------------------
% TABELA EMENTA 6: FÍSICA I
% ---------------------------------------------------------
\begin{table}[H]
\centering
\small
\begin{tabularx}{\linewidth}{|p{3.4cm}|X|p{3.5cm}|X|}
\hline
\multicolumn{4}{|>{\centering\arraybackslash\cellcolor{ifscgreen!20}}p{\dimexpr\linewidth-2\tabcolsep-2\arrayrulewidth\relax}|}{\textbf{\large Unidade Curricular: FÍSICA I}} \\ \hline
\textbf{CH Total*:} & 80 horas & \textbf{Ano / Semestre*:} & 1º Ano (Semestres 1 e 2) \\ \hline
\textbf{CH* Prática:} & 20 horas & \textbf{CH EaD*:} & não se aplica \\ \hline
\multicolumn{2}{|l|}{\textbf{CH com Divisão de Turma*:}} & \multicolumn{2}{l|}{não se aplica} \\ \hline
\multicolumn{4}{|p{\dimexpr\linewidth-2\tabcolsep-2\arrayrulewidth\relax}|}{\cellcolor{cinzaTabela}\textbf{Objetivos:}} \\ \hline
\multicolumn{4}{|p{\dimexpr\linewidth-2\tabcolsep-2\arrayrulewidth\relax}|}{
\begin{itemize}[leftmargin=*,noitemsep,topsep=2pt]
    \item Compreender os princípios da mecânica clássica e leis do movimento.
\end{itemize}
} \\ \hline
\multicolumn{4}{|p{\dimexpr\linewidth-2\tabcolsep-2\arrayrulewidth\relax}|}{\cellcolor{cinzaTabela}\textbf{Conteúdos:}} \\ \hline
\multicolumn{4}{|p{\dimexpr\linewidth-2\tabcolsep-2\arrayrulewidth\relax}|}{
\begin{itemize}[leftmargin=*,noitemsep,topsep=2pt]
    \item Cinemática Escalar e Vetorial.
    \item Leis de Newton e Aplicações.
    \item Trabalho, Energia e Conservação.
\end{itemize}
} \\ \hline
\multicolumn{4}{|p{\dimexpr\linewidth-2\tabcolsep-2\arrayrulewidth\relax}|}{\cellcolor{cinzaTabela}\textbf{Metodologia de Abordagem:}} \\ \hline
\multicolumn{4}{|p{\dimexpr\linewidth-2\tabcolsep-2\arrayrulewidth\relax}|}{Aulas teóricas e experimentos em laboratório de física. Provas e relatórios de experimentos.} \\ \hline
\multicolumn{4}{|p{\dimexpr\linewidth-2\tabcolsep-2\arrayrulewidth\relax}|}{\cellcolor{cinzaTabela}\textbf{Bibliografia Básica:}} \\ \hline
\multicolumn{4}{|p{\dimexpr\linewidth-2\tabcolsep-2\arrayrulewidth\relax}|}{1. HALLIDAY, David; RESNICK, Robert; WALKER, Jearl. Fundamentos de Física: Mecânica. 10. ed. Rio de Janeiro: LTC, 2016.} \\ \hline
\multicolumn{4}{|p{\dimexpr\linewidth-2\tabcolsep-2\arrayrulewidth\relax}|}{\cellcolor{cinzaTabela}\textbf{Bibliografia Complementar:}} \\ \hline
\multicolumn{4}{|p{\dimexpr\linewidth-2\tabcolsep-2\arrayrulewidth\relax}|}{1. TIPLER, Paul A.; MOSCA, Gene. Física para Cientistas e Engenheiros. 6. ed. Rio de Janeiro: LTC, 2014.} \\ \hline
\end{tabularx}
\end{table}

\vspace{0.3cm}


% ---------------------------------------------------------
% TABELA EMENTA 7: QUÍMICA I
% ---------------------------------------------------------
\begin{table}[H]
\centering
\small
\begin{tabularx}{\linewidth}{|p{3.4cm}|X|p{3.5cm}|X|}
\hline
\multicolumn{4}{|>{\centering\arraybackslash\cellcolor{ifscgreen!20}}p{\dimexpr\linewidth-2\tabcolsep-2\arrayrulewidth\relax}|}{\textbf{\large Unidade Curricular: QUÍMICA I}} \\ \hline
\textbf{CH Total*:} & 80 horas & \textbf{Ano / Semestre*:} & 1º Ano (Semestres 1 e 2) \\ \hline
\textbf{CH* Prática:} & 20 horas & \textbf{CH EaD*:} & não se aplica \\ \hline
\multicolumn{2}{|l|}{\textbf{CH com Divisão de Turma*:}} & \multicolumn{2}{l|}{não se aplica} \\ \hline
\multicolumn{4}{|p{\dimexpr\linewidth-2\tabcolsep-2\arrayrulewidth\relax}|}{\cellcolor{cinzaTabela}\textbf{Objetivos:}} \\ \hline
\multicolumn{4}{|p{\dimexpr\linewidth-2\tabcolsep-2\arrayrulewidth\relax}|}{
\begin{itemize}[leftmargin=*,noitemsep,topsep=2pt]
    \item Compreender a matéria, ligações químicas e tabela periódica.
\end{itemize}
} \\ \hline
\multicolumn{4}{|p{\dimexpr\linewidth-2\tabcolsep-2\arrayrulewidth\relax}|}{\cellcolor{cinzaTabela}\textbf{Conteúdos:}} \\ \hline
\multicolumn{4}{|p{\dimexpr\linewidth-2\tabcolsep-2\arrayrulewidth\relax}|}{
\begin{itemize}[leftmargin=*,noitemsep,topsep=2pt]
    \item Propriedades da Matéria e Modelos Atômicos.
    \item Tabela Periódica e Ligações Químicas.
    \item Funções Inorgânicas.
\end{itemize}
} \\ \hline
\multicolumn{4}{|p{\dimexpr\linewidth-2\tabcolsep-2\arrayrulewidth\relax}|}{\cellcolor{cinzaTabela}\textbf{Metodologia de Abordagem:}} \\ \hline
\multicolumn{4}{|p{\dimexpr\linewidth-2\tabcolsep-2\arrayrulewidth\relax}|}{Aulas dialogadas e práticas em laboratório de química. Provas escritas e trabalhos de pesquisa.} \\ \hline
\multicolumn{4}{|p{\dimexpr\linewidth-2\tabcolsep-2\arrayrulewidth\relax}|}{\cellcolor{cinzaTabela}\textbf{Bibliografia Básica:}} \\ \hline
\multicolumn{4}{|p{\dimexpr\linewidth-2\tabcolsep-2\arrayrulewidth\relax}|}{1. RUSSELL, John B. Química Geral. 2. ed. São Paulo: Makron Books, 2018.} \\ \hline
\multicolumn{4}{|p{\dimexpr\linewidth-2\tabcolsep-2\arrayrulewidth\relax}|}{\cellcolor{cinzaTabela}\textbf{Bibliografia Complementar:}} \\ \hline
\multicolumn{4}{|p{\dimexpr\linewidth-2\tabcolsep-2\arrayrulewidth\relax}|}{1. ATKINS, Peter; JONES, Loretta. Princípios de Química. 5. ed. Porto Alegre: Bookman, 2012.} \\ \hline
\end{tabularx}
\end{table}

\vspace{0.3cm}


% ---------------------------------------------------------
% TABELA EMENTA 8: FILOSOFIA I
% ---------------------------------------------------------
\begin{table}[H]
\centering
\small
\begin{tabularx}{\linewidth}{|p{3.4cm}|X|p{3.5cm}|X|}
\hline
\multicolumn{4}{|>{\centering\arraybackslash\cellcolor{ifscgreen!20}}p{\dimexpr\linewidth-2\tabcolsep-2\arrayrulewidth\relax}|}{\textbf{\large Unidade Curricular: FILOSOFIA I}} \\ \hline
\textbf{CH Total*:} & 80 horas & \textbf{Ano / Semestre*:} & 1º Ano (Semestres 1 e 2) \\ \hline
\textbf{CH* Prática:} & 0 horas & \textbf{CH EaD*:} & não se aplica \\ \hline
\multicolumn{2}{|l|}{\textbf{CH com Divisão de Turma*:}} & \multicolumn{2}{l|}{não se aplica} \\ \hline
\multicolumn{4}{|p{\dimexpr\linewidth-2\tabcolsep-2\arrayrulewidth\relax}|}{\cellcolor{cinzaTabela}\textbf{Objetivos:}} \\ \hline
\multicolumn{4}{|p{\dimexpr\linewidth-2\tabcolsep-2\arrayrulewidth\relax}|}{
\begin{itemize}[leftmargin=*,noitemsep,topsep=2pt]
    \item Desenvolver a reflexão filosófica crítica sobre ética, política e conhecimento.
\end{itemize}
} \\ \hline
\multicolumn{4}{|p{\dimexpr\linewidth-2\tabcolsep-2\arrayrulewidth\relax}|}{\cellcolor{cinzaTabela}\textbf{Conteúdos:}} \\ \hline
\multicolumn{4}{|p{\dimexpr\linewidth-2\tabcolsep-2\arrayrulewidth\relax}|}{
\begin{itemize}[leftmargin=*,noitemsep,topsep=2pt]
    \item Origens do Pensamento Filosófico.
    \item Ética e Filosofia Política Antiga.
    \item Teoria do Conhecimento.
\end{itemize}
} \\ \hline
\multicolumn{4}{|p{\dimexpr\linewidth-2\tabcolsep-2\arrayrulewidth\relax}|}{\cellcolor{cinzaTabela}\textbf{Metodologia de Abordagem:}} \\ \hline
\multicolumn{4}{|p{\dimexpr\linewidth-2\tabcolsep-2\arrayrulewidth\relax}|}{Leitura guiada de clássicos e debates em sala. Produção textual reflexiva e participação.} \\ \hline
\multicolumn{4}{|p{\dimexpr\linewidth-2\tabcolsep-2\arrayrulewidth\relax}|}{\cellcolor{cinzaTabela}\textbf{Bibliografia Básica:}} \\ \hline
\multicolumn{4}{|p{\dimexpr\linewidth-2\tabcolsep-2\arrayrulewidth\relax}|}{1. CHAUÍ, Marilena. Convite à Filosofia. 14. ed. São Paulo: Ática, 2020.} \\ \hline
\multicolumn{4}{|p{\dimexpr\linewidth-2\tabcolsep-2\arrayrulewidth\relax}|}{\cellcolor{cinzaTabela}\textbf{Bibliografia Complementar:}} \\ \hline
\multicolumn{4}{|p{\dimexpr\linewidth-2\tabcolsep-2\arrayrulewidth\relax}|}{1. ARANHA, Maria Lúcia de Arruda. Filosofando. São Paulo: Moderna, 2019.} \\ \hline
\end{tabularx}
\end{table}

\vspace{0.3cm}


% ---------------------------------------------------------
% TABELA EMENTA 9: EDUCAÇÃO FÍSICA I
% ---------------------------------------------------------
\begin{table}[H]
\centering
\small
\begin{tabularx}{\linewidth}{|p{3.4cm}|X|p{3.5cm}|X|}
\hline
\multicolumn{4}{|>{\centering\arraybackslash\cellcolor{ifscgreen!20}}p{\dimexpr\linewidth-2\tabcolsep-2\arrayrulewidth\relax}|}{\textbf{\large Unidade Curricular: EDUCAÇÃO FÍSICA I}} \\ \hline
\textbf{CH Total*:} & 80 horas & \textbf{Ano / Semestre*:} & 1º Ano (Semestres 1 e 2) \\ \hline
\textbf{CH* Prática:} & 60 horas & \textbf{CH EaD*:} & não se aplica \\ \hline
\multicolumn{2}{|l|}{\textbf{CH com Divisão de Turma*:}} & \multicolumn{2}{l|}{não se aplica} \\ \hline
\multicolumn{4}{|p{\dimexpr\linewidth-2\tabcolsep-2\arrayrulewidth\relax}|}{\cellcolor{cinzaTabela}\textbf{Objetivos:}} \\ \hline
\multicolumn{4}{|p{\dimexpr\linewidth-2\tabcolsep-2\arrayrulewidth\relax}|}{
\begin{itemize}[leftmargin=*,noitemsep,topsep=2pt]
    \item Promover a saúde, qualidade de vida e vivência de práticas corporais coletivas.
\end{itemize}
} \\ \hline
\multicolumn{4}{|p{\dimexpr\linewidth-2\tabcolsep-2\arrayrulewidth\relax}|}{\cellcolor{cinzaTabela}\textbf{Conteúdos:}} \\ \hline
\multicolumn{4}{|p{\dimexpr\linewidth-2\tabcolsep-2\arrayrulewidth\relax}|}{
\begin{itemize}[leftmargin=*,noitemsep,topsep=2pt]
    \item Cultura Corporal e Saúde.
    \item Esportes Coletivos e Jogos.
    \item Ergonomia Corporal e Postura.
\end{itemize}
} \\ \hline
\multicolumn{4}{|p{\dimexpr\linewidth-2\tabcolsep-2\arrayrulewidth\relax}|}{\cellcolor{cinzaTabela}\textbf{Metodologia de Abordagem:}} \\ \hline
\multicolumn{4}{|p{\dimexpr\linewidth-2\tabcolsep-2\arrayrulewidth\relax}|}{Práticas corporais no ginásio/quadra e aulas teóricas sobre saúde. Avaliação pela participação e vivências.} \\ \hline
\multicolumn{4}{|p{\dimexpr\linewidth-2\tabcolsep-2\arrayrulewidth\relax}|}{\cellcolor{cinzaTabela}\textbf{Bibliografia Básica:}} \\ \hline
\multicolumn{4}{|p{\dimexpr\linewidth-2\tabcolsep-2\arrayrulewidth\relax}|}{1. DARIDO, Suraya Cristina. Educação Física na Escola. Rio de Janeiro: Guanabara Koogan, 2017.} \\ \hline
\multicolumn{4}{|p{\dimexpr\linewidth-2\tabcolsep-2\arrayrulewidth\relax}|}{\cellcolor{cinzaTabela}\textbf{Bibliografia Complementar:}} \\ \hline
\multicolumn{4}{|p{\dimexpr\linewidth-2\tabcolsep-2\arrayrulewidth\relax}|}{1. KUNZ, Elenor. Transformação Didático-Pedagógica do Esporte. Ijuí: Unijuí, 2014.} \\ \hline
\end{tabularx}
\end{table}

\vspace{0.3cm}


% ---------------------------------------------------------
% TABELA EMENTA 10: ESPANHOL
% ---------------------------------------------------------
\begin{table}[H]
\centering
\small
\begin{tabularx}{\linewidth}{|p{3.4cm}|X|p{3.5cm}|X|}
\hline
\multicolumn{4}{|>{\centering\arraybackslash\cellcolor{ifscgreen!20}}p{\dimexpr\linewidth-2\tabcolsep-2\arrayrulewidth\relax}|}{\textbf{\large Unidade Curricular: ESPANHOL}} \\ \hline
\textbf{CH Total*:} & 80 horas & \textbf{Ano / Semestre*:} & 1º Ano (Semestres 1 e 2) \\ \hline
\textbf{CH* Prática:} & 20 horas & \textbf{CH EaD*:} & não se aplica \\ \hline
\multicolumn{2}{|l|}{\textbf{CH com Divisão de Turma*:}} & \multicolumn{2}{l|}{não se aplica} \\ \hline
\multicolumn{4}{|p{\dimexpr\linewidth-2\tabcolsep-2\arrayrulewidth\relax}|}{\cellcolor{cinzaTabela}\textbf{Objetivos:}} \\ \hline
\multicolumn{4}{|p{\dimexpr\linewidth-2\tabcolsep-2\arrayrulewidth\relax}|}{
\begin{itemize}[leftmargin=*,noitemsep,topsep=2pt]
    \item Compreender e comunicar-se em língua espanhola em situações cotidianas e profissionais.
\end{itemize}
} \\ \hline
\multicolumn{4}{|p{\dimexpr\linewidth-2\tabcolsep-2\arrayrulewidth\relax}|}{\cellcolor{cinzaTabela}\textbf{Conteúdos:}} \\ \hline
\multicolumn{4}{|p{\dimexpr\linewidth-2\tabcolsep-2\arrayrulewidth\relax}|}{
\begin{itemize}[leftmargin=*,noitemsep,topsep=2pt]
    \item Vocabulário Básico e Gramática Fundamental.
    \item Leitura e Interpretação de Textos.
    \item Comunicação Oral e Escrita em Espanhol.
\end{itemize}
} \\ \hline
\multicolumn{4}{|p{\dimexpr\linewidth-2\tabcolsep-2\arrayrulewidth\relax}|}{\cellcolor{cinzaTabela}\textbf{Metodologia de Abordagem:}} \\ \hline
\multicolumn{4}{|p{\dimexpr\linewidth-2\tabcolsep-2\arrayrulewidth\relax}|}{Aulas práticas interativas de conversação e interpretação textual. Avaliações orais e escritas.} \\ \hline
\multicolumn{4}{|p{\dimexpr\linewidth-2\tabcolsep-2\arrayrulewidth\relax}|}{\cellcolor{cinzaTabela}\textbf{Bibliografia Básica:}} \\ \hline
\multicolumn{4}{|p{\dimexpr\linewidth-2\tabcolsep-2\arrayrulewidth\relax}|}{1. COIMBRA, Ludmila et al. Cercanías: Español para Brasileños. São Paulo: SM, 2018.} \\ \hline
\multicolumn{4}{|p{\dimexpr\linewidth-2\tabcolsep-2\arrayrulewidth\relax}|}{\cellcolor{cinzaTabela}\textbf{Bibliografia Complementar:}} \\ \hline
\multicolumn{4}{|p{\dimexpr\linewidth-2\tabcolsep-2\arrayrulewidth\relax}|}{1. FANJUL, Adrián. Gramática de Español para Hispanohablantes. São Paulo: Santillana, 2014.} \\ \hline
\end{tabularx}
\end{table}

\vspace{0.3cm}


% ---------------------------------------------------------
% TABELA EMENTA 11: INTRODUÇÃO À ADMINISTRAÇÃO
% ---------------------------------------------------------
\begin{table}[H]
\centering
\small
\begin{tabularx}{\linewidth}{|p{3.4cm}|X|p{3.5cm}|X|}
\hline
\multicolumn{4}{|>{\centering\arraybackslash\cellcolor{ifscgreen!20}}p{\dimexpr\linewidth-2\tabcolsep-2\arrayrulewidth\relax}|}{\textbf{\large Unidade Curricular: INTRODUÇÃO À ADMINISTRAÇÃO}} \\ \hline
\textbf{CH Total*:} & 80 horas & \textbf{Ano / Semestre*:} & 1º Ano (Semestre 1) \\ \hline
\textbf{CH* Prática:} & 30 horas & \textbf{CH EaD*:} & não se aplica \\ \hline
\multicolumn{2}{|l|}{\textbf{CH com Divisão de Turma*:}} & \multicolumn{2}{l|}{não se aplica} \\ \hline
\multicolumn{4}{|p{\dimexpr\linewidth-2\tabcolsep-2\arrayrulewidth\relax}|}{\cellcolor{cinzaTabela}\textbf{Objetivos:}} \\ \hline
\multicolumn{4}{|p{\dimexpr\linewidth-2\tabcolsep-2\arrayrulewidth\relax}|}{
\begin{itemize}[leftmargin=*,noitemsep,topsep=2pt]
    \item Compreender a evolução histórica das teorias administrativas e suas aplicações modernas nas empresas;
    \item Identificar as funções fundamentais de planejamento, organização, direção e controle (PODC);
    \item Desenvolver visão analítica sobre ética, governança e responsabilidade socioambiental corporativa.
\end{itemize}
} \\ \hline
\multicolumn{4}{|p{\dimexpr\linewidth-2\tabcolsep-2\arrayrulewidth\relax}|}{\cellcolor{cinzaTabela}\textbf{Conteúdos:}} \\ \hline
\multicolumn{4}{|p{\dimexpr\linewidth-2\tabcolsep-2\arrayrulewidth\relax}|}{
\begin{itemize}[leftmargin=*,noitemsep,topsep=2pt]
    \item Evolução do Pensamento Administrativo e Teorias Organizacionais.
    \item As Funções Administrativas: Planejamento Estratégico, Tático e Operacional.
    \item Organização, Direção (Liderança e Motivação) e Controle.
    \item Ética, Governança e Sustentabilidade Organizacional.
\end{itemize}
} \\ \hline
\multicolumn{4}{|p{\dimexpr\linewidth-2\tabcolsep-2\arrayrulewidth\relax}|}{\cellcolor{cinzaTabela}\textbf{Metodologia de Abordagem:}} \\ \hline
\multicolumn{4}{|p{\dimexpr\linewidth-2\tabcolsep-2\arrayrulewidth\relax}|}{Aulas expositivo-dialogadas, análise de casos práticos e seminários corporativos. Avaliação por estudos de caso e provas teóricas.} \\ \hline
\multicolumn{4}{|p{\dimexpr\linewidth-2\tabcolsep-2\arrayrulewidth\relax}|}{\cellcolor{cinzaTabela}\textbf{Bibliografia Básica:}} \\ \hline
\multicolumn{4}{|p{\dimexpr\linewidth-2\tabcolsep-2\arrayrulewidth\relax}|}{1. CHIAVENATO, Idalberto. Introdução à Teoria Geral da Administração. 10. ed. Rio de Janeiro: GEN LTC, 2021.
2. MAXIMIANO, Antonio Cesar Amaru. Introdução à Administração. 8. ed. São Paulo: Atlas, 2017.} \\ \hline
\multicolumn{4}{|p{\dimexpr\linewidth-2\tabcolsep-2\arrayrulewidth\relax}|}{\cellcolor{cinzaTabela}\textbf{Bibliografia Complementar:}} \\ \hline
\multicolumn{4}{|p{\dimexpr\linewidth-2\tabcolsep-2\arrayrulewidth\relax}|}{1. SOBRAL, Filipe; PECI, Alketa. Administração: teoria e prática no contexto brasileiro. 2. ed. São Paulo: Pearson, 2013.} \\ \hline
\end{tabularx}
\end{table}

\vspace{0.3cm}


% ---------------------------------------------------------
% TABELA EMENTA 12: INFORMÁTICA APLICADA
% ---------------------------------------------------------
\begin{table}[H]
\centering
\small
\begin{tabularx}{\linewidth}{|p{3.4cm}|X|p{3.5cm}|X|}
\hline
\multicolumn{4}{|>{\centering\arraybackslash\cellcolor{ifscgreen!20}}p{\dimexpr\linewidth-2\tabcolsep-2\arrayrulewidth\relax}|}{\textbf{\large Unidade Curricular: INFORMÁTICA APLICADA}} \\ \hline
\textbf{CH Total*:} & 60 horas & \textbf{Ano / Semestre*:} & 1º Ano (Semestre 1) \\ \hline
\textbf{CH* Prática:} & 40 horas & \textbf{CH EaD*:} & não se aplica \\ \hline
\multicolumn{2}{|l|}{\textbf{CH com Divisão de Turma*:}} & \multicolumn{2}{l|}{não se aplica} \\ \hline
\multicolumn{4}{|p{\dimexpr\linewidth-2\tabcolsep-2\arrayrulewidth\relax}|}{\cellcolor{cinzaTabela}\textbf{Objetivos:}} \\ \hline
\multicolumn{4}{|p{\dimexpr\linewidth-2\tabcolsep-2\arrayrulewidth\relax}|}{
\begin{itemize}[leftmargin=*,noitemsep,topsep=2pt]
    \item Dominar recursos essenciais de edição de documentos corporativos e planilhas eletrônicas;
    \item Aplicar ferramentas computacionais na resolução de problemas administrativos e financeiros;
\end{itemize}
} \\ \hline
\multicolumn{4}{|p{\dimexpr\linewidth-2\tabcolsep-2\arrayrulewidth\relax}|}{\cellcolor{cinzaTabela}\textbf{Conteúdos:}} \\ \hline
\multicolumn{4}{|p{\dimexpr\linewidth-2\tabcolsep-2\arrayrulewidth\relax}|}{
\begin{itemize}[leftmargin=*,noitemsep,topsep=2pt]
    \item Sistemas Operacionais, Arquivos em Nuvem e Segurança da Informação.
    \item Processamento de Textos Corporativos e Formatação Normativa.
    \item Planilhas Eletrônicas para Gestão (Fórmulas, Funções Lógicas e Tabelas Dinâmicas).
\end{itemize}
} \\ \hline
\multicolumn{4}{|p{\dimexpr\linewidth-2\tabcolsep-2\arrayrulewidth\relax}|}{\cellcolor{cinzaTabela}\textbf{Metodologia de Abordagem:}} \\ \hline
\multicolumn{4}{|p{\dimexpr\linewidth-2\tabcolsep-2\arrayrulewidth\relax}|}{Aulas práticas em laboratório de informática com resolução de exercícios orientados. Avaliação prática no computador.} \\ \hline
\multicolumn{4}{|p{\dimexpr\linewidth-2\tabcolsep-2\arrayrulewidth\relax}|}{\cellcolor{cinzaTabela}\textbf{Bibliografia Básica:}} \\ \hline
\multicolumn{4}{|p{\dimexpr\linewidth-2\tabcolsep-2\arrayrulewidth\relax}|}{1. MANZANO, José Augusto N. G. Estudo Dirigido de Microsoft Office Excel 2019. São Paulo: Érica, 2019.
2. TURBAN, Efraim et al. Tecnologia da Informação para Gestão. 8. ed. Porto Alegre: Bookman, 2013.} \\ \hline
\multicolumn{4}{|p{\dimexpr\linewidth-2\tabcolsep-2\arrayrulewidth\relax}|}{\cellcolor{cinzaTabela}\textbf{Bibliografia Complementar:}} \\ \hline
\multicolumn{4}{|p{\dimexpr\linewidth-2\tabcolsep-2\arrayrulewidth\relax}|}{1. STAIR, Ralph M.; REYNOLDS, George W. Princípios de Sistemas de Informação. 12. ed. Cengage, 2018.} \\ \hline
\end{tabularx}
\end{table}

\vspace{0.3cm}


% ---------------------------------------------------------
% TABELA EMENTA 13: ORGANIZAÇÃO E PROCESSOS
% ---------------------------------------------------------
\begin{table}[H]
\centering
\small
\begin{tabularx}{\linewidth}{|p{3.4cm}|X|p{3.5cm}|X|}
\hline
\multicolumn{4}{|>{\centering\arraybackslash\cellcolor{ifscgreen!20}}p{\dimexpr\linewidth-2\tabcolsep-2\arrayrulewidth\relax}|}{\textbf{\large Unidade Curricular: ORGANIZAÇÃO E PROCESSOS}} \\ \hline
\textbf{CH Total*:} & 80 horas & \textbf{Ano / Semestre*:} & 1º Ano (Semestre 2) \\ \hline
\textbf{CH* Prática:} & 30 horas & \textbf{CH EaD*:} & não se aplica \\ \hline
\multicolumn{2}{|l|}{\textbf{CH com Divisão de Turma*:}} & \multicolumn{2}{l|}{não se aplica} \\ \hline
\multicolumn{4}{|p{\dimexpr\linewidth-2\tabcolsep-2\arrayrulewidth\relax}|}{\cellcolor{cinzaTabela}\textbf{Objetivos:}} \\ \hline
\multicolumn{4}{|p{\dimexpr\linewidth-2\tabcolsep-2\arrayrulewidth\relax}|}{
\begin{itemize}[leftmargin=*,noitemsep,topsep=2pt]
    \item Mapear e desenhar fluxogramas de processos organizacionais (BPM);
    \item Propor melhorias no arranjo físico e na ergonomia dos postos de trabalho;
\end{itemize}
} \\ \hline
\multicolumn{4}{|p{\dimexpr\linewidth-2\tabcolsep-2\arrayrulewidth\relax}|}{\cellcolor{cinzaTabela}\textbf{Conteúdos:}} \\ \hline
\multicolumn{4}{|p{\dimexpr\linewidth-2\tabcolsep-2\arrayrulewidth\relax}|}{
\begin{itemize}[leftmargin=*,noitemsep,topsep=2pt]
    \item Estruturas Organizacionais (Funcional, Divisional, Matricial).
    \item Gestão por Processos (BPM) e Mapeamento com Fluxogramas.
    \item Elaboração de Manuais Administrativos e Formulários.
    \item Estudo de Tempos, Métodos, Layout e Ergonomia.
\end{itemize}
} \\ \hline
\multicolumn{4}{|p{\dimexpr\linewidth-2\tabcolsep-2\arrayrulewidth\relax}|}{\cellcolor{cinzaTabela}\textbf{Metodologia de Abordagem:}} \\ \hline
\multicolumn{4}{|p{\dimexpr\linewidth-2\tabcolsep-2\arrayrulewidth\relax}|}{Aulas práticas de mapeamento de processos e desenho de fluxogramas em laboratório. Avaliação por elaboração de manual de processos.} \\ \hline
\multicolumn{4}{|p{\dimexpr\linewidth-2\tabcolsep-2\arrayrulewidth\relax}|}{\cellcolor{cinzaTabela}\textbf{Bibliografia Básica:}} \\ \hline
\multicolumn{4}{|p{\dimexpr\linewidth-2\tabcolsep-2\arrayrulewidth\relax}|}{1. CURURY, Antonio. Organização e Métodos: uma visão holística. 8. ed. São Paulo: Atlas, 2018.
2. GONÇALVES, José Ernesto Lima. Gestão por Processos na Prática. São Paulo: RAE, 2020.} \\ \hline
\multicolumn{4}{|p{\dimexpr\linewidth-2\tabcolsep-2\arrayrulewidth\relax}|}{\cellcolor{cinzaTabela}\textbf{Bibliografia Complementar:}} \\ \hline
\multicolumn{4}{|p{\dimexpr\linewidth-2\tabcolsep-2\arrayrulewidth\relax}|}{1. OLIVEIRA, Djalma de Pinho Rebouças de. Sistemas, Organização e Métodos. 21. ed. São Paulo: Atlas, 2013.} \\ \hline
\end{tabularx}
\end{table}

\vspace{0.3cm}


% ---------------------------------------------------------
% TABELA EMENTA 14: SOCIEDADE E TRABALHO
% ---------------------------------------------------------
\begin{table}[H]
\centering
\small
\begin{tabularx}{\linewidth}{|p{3.4cm}|X|p{3.5cm}|X|}
\hline
\multicolumn{4}{|>{\centering\arraybackslash\cellcolor{ifscgreen!20}}p{\dimexpr\linewidth-2\tabcolsep-2\arrayrulewidth\relax}|}{\textbf{\large Unidade Curricular: SOCIEDADE E TRABALHO}} \\ \hline
\textbf{CH Total*:} & 60 horas & \textbf{Ano / Semestre*:} & 1º Ano (Semestre 2) \\ \hline
\textbf{CH* Prática:} & 15 horas & \textbf{CH EaD*:} & não se aplica \\ \hline
\multicolumn{2}{|l|}{\textbf{CH com Divisão de Turma*:}} & \multicolumn{2}{l|}{não se aplica} \\ \hline
\multicolumn{4}{|p{\dimexpr\linewidth-2\tabcolsep-2\arrayrulewidth\relax}|}{\cellcolor{cinzaTabela}\textbf{Objetivos:}} \\ \hline
\multicolumn{4}{|p{\dimexpr\linewidth-2\tabcolsep-2\arrayrulewidth\relax}|}{
\begin{itemize}[leftmargin=*,noitemsep,topsep=2pt]
    \item Analisar as transformações históricas do mundo do trabalho e seus impactos sociais;
    \item Compreender a cidadania e os direitos humanos como pilares da atuação profissional ética.
\end{itemize}
} \\ \hline
\multicolumn{4}{|p{\dimexpr\linewidth-2\tabcolsep-2\arrayrulewidth\relax}|}{\cellcolor{cinzaTabela}\textbf{Conteúdos:}} \\ \hline
\multicolumn{4}{|p{\dimexpr\linewidth-2\tabcolsep-2\arrayrulewidth\relax}|}{
\begin{itemize}[leftmargin=*,noitemsep,topsep=2pt]
    \item História do Trabalho e Mudanças no Mundo Contemporâneo.
    \item Cidadania, Direitos Humanos e Inclusão Social nas Empresas.
    \item Ética Profissional e Responsabilidade Social Corporativa.
\end{itemize}
} \\ \hline
\multicolumn{4}{|p{\dimexpr\linewidth-2\tabcolsep-2\arrayrulewidth\relax}|}{\cellcolor{cinzaTabela}\textbf{Metodologia de Abordagem:}} \\ \hline
\multicolumn{4}{|p{\dimexpr\linewidth-2\tabcolsep-2\arrayrulewidth\relax}|}{Debates, exibição de documentários e análise crítica de textos. Avaliação por produção textual dissertativa e seminários.} \\ \hline
\multicolumn{4}{|p{\dimexpr\linewidth-2\tabcolsep-2\arrayrulewidth\relax}|}{\cellcolor{cinzaTabela}\textbf{Bibliografia Básica:}} \\ \hline
\multicolumn{4}{|p{\dimexpr\linewidth-2\tabcolsep-2\arrayrulewidth\relax}|}{1. ANTUNES, Ricardo. Os Sentidos do Trabalho. São Paulo: Boitempo, 2009.} \\ \hline
\multicolumn{4}{|p{\dimexpr\linewidth-2\tabcolsep-2\arrayrulewidth\relax}|}{\cellcolor{cinzaTabela}\textbf{Bibliografia Complementar:}} \\ \hline
\multicolumn{4}{|p{\dimexpr\linewidth-2\tabcolsep-2\arrayrulewidth\relax}|}{1. SINGELMANN, Joachim. Transformações do Trabalho na Sociedade Moderna. São Paulo: Ed. 34, 2015.} \\ \hline
\end{tabularx}
\end{table}

\vspace{0.3cm}


% ---------------------------------------------------------
% TABELA EMENTA 15: OFICINA DE INTEGRAÇÃO I
% ---------------------------------------------------------
\begin{table}[H]
\centering
\small
\begin{tabularx}{\linewidth}{|p{3.4cm}|X|p{3.5cm}|X|}
\hline
\multicolumn{4}{|>{\centering\arraybackslash\cellcolor{ifscgreen!20}}p{\dimexpr\linewidth-2\tabcolsep-2\arrayrulewidth\relax}|}{\textbf{\large Unidade Curricular: OFICINA DE INTEGRAÇÃO I}} \\ \hline
\textbf{CH Total*:} & 40 horas & \textbf{Ano / Semestre*:} & 1º Ano (Semestres 1 e 2) \\ \hline
\textbf{CH* Prática:} & 30 horas & \textbf{CH EaD*:} & não se aplica \\ \hline
\multicolumn{2}{|l|}{\textbf{CH com Divisão de Turma*:}} & \multicolumn{2}{l|}{não se aplica} \\ \hline
\multicolumn{4}{|p{\dimexpr\linewidth-2\tabcolsep-2\arrayrulewidth\relax}|}{\cellcolor{cinzaTabela}\textbf{Objetivos:}} \\ \hline
\multicolumn{4}{|p{\dimexpr\linewidth-2\tabcolsep-2\arrayrulewidth\relax}|}{
\begin{itemize}[leftmargin=*,noitemsep,topsep=2pt]
    \item Promover a integração discente e o alinhamento de conhecimentos fundamentais;
    \item Desenvolver competências de organização do estudo e métodos de aprendizagem.
\end{itemize}
} \\ \hline
\multicolumn{4}{|p{\dimexpr\linewidth-2\tabcolsep-2\arrayrulewidth\relax}|}{\cellcolor{cinzaTabela}\textbf{Conteúdos:}} \\ \hline
\multicolumn{4}{|p{\dimexpr\linewidth-2\tabcolsep-2\arrayrulewidth\relax}|}{
\begin{itemize}[leftmargin=*,noitemsep,topsep=2pt]
    \item Métodos de Estudo e Organização do Tempo Acadêmico.
    \item Leitura Crítica, Redação e Pesquisa Bibliográfica Inicial.
    \item Integração entre os Blocos Geral e Técnico do Curso.
\end{itemize}
} \\ \hline
\multicolumn{4}{|p{\dimexpr\linewidth-2\tabcolsep-2\arrayrulewidth\relax}|}{\cellcolor{cinzaTabela}\textbf{Metodologia de Abordagem:}} \\ \hline
\multicolumn{4}{|p{\dimexpr\linewidth-2\tabcolsep-2\arrayrulewidth\relax}|}{Oficinas práticas, dinâmicas de grupo e seminários de acolhimento. Avaliação por relatórios de participação.} \\ \hline
\multicolumn{4}{|p{\dimexpr\linewidth-2\tabcolsep-2\arrayrulewidth\relax}|}{\cellcolor{cinzaTabela}\textbf{Bibliografia Básica:}} \\ \hline
\multicolumn{4}{|p{\dimexpr\linewidth-2\tabcolsep-2\arrayrulewidth\relax}|}{1. SEVERINO, Antônio Joaquim. Metodologia do Trabalho Científico. 24. ed. São Paulo: Cortez, 2017.} \\ \hline
\multicolumn{4}{|p{\dimexpr\linewidth-2\tabcolsep-2\arrayrulewidth\relax}|}{\cellcolor{cinzaTabela}\textbf{Bibliografia Complementar:}} \\ \hline
\multicolumn{4}{|p{\dimexpr\linewidth-2\tabcolsep-2\arrayrulewidth\relax}|}{1. GIL, Antonio Carlos. Como Elaborar Projetos de Pesquisa. 6. ed. São Paulo: Atlas, 2019.} \\ \hline
\end{tabularx}
\end{table}

\vspace{0.3cm}


% ---------------------------------------------------------
% TABELA EMENTA 16: LÍNGUA PORTUGUESA E LITERATURA II
% ---------------------------------------------------------
\begin{table}[H]
\centering
\small
\begin{tabularx}{\linewidth}{|p{3.4cm}|X|p{3.5cm}|X|}
\hline
\multicolumn{4}{|>{\centering\arraybackslash\cellcolor{ifscgreen!20}}p{\dimexpr\linewidth-2\tabcolsep-2\arrayrulewidth\relax}|}{\textbf{\large Unidade Curricular: LÍNGUA PORTUGUESA E LITERATURA II}} \\ \hline
\textbf{CH Total*:} & 80 horas & \textbf{Ano / Semestre*:} & 2º Ano (Semestres 1 e 2) \\ \hline
\textbf{CH* Prática:} & 20 horas & \textbf{CH EaD*:} & não se aplica \\ \hline
\multicolumn{2}{|l|}{\textbf{CH com Divisão de Turma*:}} & \multicolumn{2}{l|}{não se aplica} \\ \hline
\multicolumn{4}{|p{\dimexpr\linewidth-2\tabcolsep-2\arrayrulewidth\relax}|}{\cellcolor{cinzaTabela}\textbf{Objetivos:}} \\ \hline
\multicolumn{4}{|p{\dimexpr\linewidth-2\tabcolsep-2\arrayrulewidth\relax}|}{
\begin{itemize}[leftmargin=*,noitemsep,topsep=2pt]
    \item Aprofundar a análise literária e a argumentação em textos acadêmicos.
\end{itemize}
} \\ \hline
\multicolumn{4}{|p{\dimexpr\linewidth-2\tabcolsep-2\arrayrulewidth\relax}|}{\cellcolor{cinzaTabela}\textbf{Conteúdos:}} \\ \hline
\multicolumn{4}{|p{\dimexpr\linewidth-2\tabcolsep-2\arrayrulewidth\relax}|}{
\begin{itemize}[leftmargin=*,noitemsep,topsep=2pt]
    \item Literatura Brasileira e Portuguesa Contemporânea.
    \item Argumentação e Coesão Textual Avançada.
    \item Redação Técnica e Corporativa.
\end{itemize}
} \\ \hline
\multicolumn{4}{|p{\dimexpr\linewidth-2\tabcolsep-2\arrayrulewidth\relax}|}{\cellcolor{cinzaTabela}\textbf{Metodologia de Abordagem:}} \\ \hline
\multicolumn{4}{|p{\dimexpr\linewidth-2\tabcolsep-2\arrayrulewidth\relax}|}{Aulas expositivas e seminários literários. Avaliação por redações e provas.} \\ \hline
\multicolumn{4}{|p{\dimexpr\linewidth-2\tabcolsep-2\arrayrulewidth\relax}|}{\cellcolor{cinzaTabela}\textbf{Bibliografia Básica:}} \\ \hline
\multicolumn{4}{|p{\dimexpr\linewidth-2\tabcolsep-2\arrayrulewidth\relax}|}{1. CEREJA, William Roberto. Literatura Brasileira. São Paulo: Atual, 2018.} \\ \hline
\multicolumn{4}{|p{\dimexpr\linewidth-2\tabcolsep-2\arrayrulewidth\relax}|}{\cellcolor{cinzaTabela}\textbf{Bibliografia Complementar:}} \\ \hline
\multicolumn{4}{|p{\dimexpr\linewidth-2\tabcolsep-2\arrayrulewidth\relax}|}{1. CANDIDO, Antonio. Formação da Literatura Brasileira. São Paulo: Ouro sobre Azul, 2014.} \\ \hline
\end{tabularx}
\end{table}

\vspace{0.3cm}


% ---------------------------------------------------------
% TABELA EMENTA 17: MATEMÁTICA II
% ---------------------------------------------------------
\begin{table}[H]
\centering
\small
\begin{tabularx}{\linewidth}{|p{3.4cm}|X|p{3.5cm}|X|}
\hline
\multicolumn{4}{|>{\centering\arraybackslash\cellcolor{ifscgreen!20}}p{\dimexpr\linewidth-2\tabcolsep-2\arrayrulewidth\relax}|}{\textbf{\large Unidade Curricular: MATEMÁTICA II}} \\ \hline
\textbf{CH Total*:} & 80 horas & \textbf{Ano / Semestre*:} & 2º Ano (Semestres 1 e 2) \\ \hline
\textbf{CH* Prática:} & 20 horas & \textbf{CH EaD*:} & não se aplica \\ \hline
\multicolumn{2}{|l|}{\textbf{CH com Divisão de Turma*:}} & \multicolumn{2}{l|}{não se aplica} \\ \hline
\multicolumn{4}{|p{\dimexpr\linewidth-2\tabcolsep-2\arrayrulewidth\relax}|}{\cellcolor{cinzaTabela}\textbf{Objetivos:}} \\ \hline
\multicolumn{4}{|p{\dimexpr\linewidth-2\tabcolsep-2\arrayrulewidth\relax}|}{
\begin{itemize}[leftmargin=*,noitemsep,topsep=2pt]
    \item Compreender matrizes, sistemas lineares e introdução à matemática financeira.
\end{itemize}
} \\ \hline
\multicolumn{4}{|p{\dimexpr\linewidth-2\tabcolsep-2\arrayrulewidth\relax}|}{\cellcolor{cinzaTabela}\textbf{Conteúdos:}} \\ \hline
\multicolumn{4}{|p{\dimexpr\linewidth-2\tabcolsep-2\arrayrulewidth\relax}|}{
\begin{itemize}[leftmargin=*,noitemsep,topsep=2pt]
    \item Matrizes, Determinantes e Sistemas Lineares.
    \item Juros Simples, Compostos e Descontos.
    \item Progressões Aritméticas e Geométricas.
\end{itemize}
} \\ \hline
\multicolumn{4}{|p{\dimexpr\linewidth-2\tabcolsep-2\arrayrulewidth\relax}|}{\cellcolor{cinzaTabela}\textbf{Metodologia de Abordagem:}} \\ \hline
\multicolumn{4}{|p{\dimexpr\linewidth-2\tabcolsep-2\arrayrulewidth\relax}|}{Aulas expositivas e resolução de problemas financeiros aplicados. Provas escritas.} \\ \hline
\multicolumn{4}{|p{\dimexpr\linewidth-2\tabcolsep-2\arrayrulewidth\relax}|}{\cellcolor{cinzaTabela}\textbf{Bibliografia Básica:}} \\ \hline
\multicolumn{4}{|p{\dimexpr\linewidth-2\tabcolsep-2\arrayrulewidth\relax}|}{1. HAZAZZ, Samer. Matemática Financeira. São Paulo: Atlas, 2019.} \\ \hline
\multicolumn{4}{|p{\dimexpr\linewidth-2\tabcolsep-2\arrayrulewidth\relax}|}{\cellcolor{cinzaTabela}\textbf{Bibliografia Complementar:}} \\ \hline
\multicolumn{4}{|p{\dimexpr\linewidth-2\tabcolsep-2\arrayrulewidth\relax}|}{1. DANTE, Luiz Roberto. Matemática. São Paulo: Ática, 2020.} \\ \hline
\end{tabularx}
\end{table}

\vspace{0.3cm}


% ---------------------------------------------------------
% TABELA EMENTA 18: HISTÓRIA II
% ---------------------------------------------------------
\begin{table}[H]
\centering
\small
\begin{tabularx}{\linewidth}{|p{3.4cm}|X|p{3.5cm}|X|}
\hline
\multicolumn{4}{|>{\centering\arraybackslash\cellcolor{ifscgreen!20}}p{\dimexpr\linewidth-2\tabcolsep-2\arrayrulewidth\relax}|}{\textbf{\large Unidade Curricular: HISTÓRIA II}} \\ \hline
\textbf{CH Total*:} & 40 horas & \textbf{Ano / Semestre*:} & 2º Ano (Semestres 1 e 2) \\ \hline
\textbf{CH* Prática:} & 0 horas & \textbf{CH EaD*:} & não se aplica \\ \hline
\multicolumn{2}{|l|}{\textbf{CH com Divisão de Turma*:}} & \multicolumn{2}{l|}{não se aplica} \\ \hline
\multicolumn{4}{|p{\dimexpr\linewidth-2\tabcolsep-2\arrayrulewidth\relax}|}{\cellcolor{cinzaTabela}\textbf{Objetivos:}} \\ \hline
\multicolumn{4}{|p{\dimexpr\linewidth-2\tabcolsep-2\arrayrulewidth\relax}|}{
\begin{itemize}[leftmargin=*,noitemsep,topsep=2pt]
    \item Compreender a história moderna e contemporânea mundial e brasileira.
\end{itemize}
} \\ \hline
\multicolumn{4}{|p{\dimexpr\linewidth-2\tabcolsep-2\arrayrulewidth\relax}|}{\cellcolor{cinzaTabela}\textbf{Conteúdos:}} \\ \hline
\multicolumn{4}{|p{\dimexpr\linewidth-2\tabcolsep-2\arrayrulewidth\relax}|}{
\begin{itemize}[leftmargin=*,noitemsep,topsep=2pt]
    \item Revolução Industrial e Iluminismo.
    \item Brasil Império e Primeira República.
    \item Conflitos Mundiais do Século XX.
\end{itemize}
} \\ \hline
\multicolumn{4}{|p{\dimexpr\linewidth-2\tabcolsep-2\arrayrulewidth\relax}|}{\cellcolor{cinzaTabela}\textbf{Metodologia de Abordagem:}} \\ \hline
\multicolumn{4}{|p{\dimexpr\linewidth-2\tabcolsep-2\arrayrulewidth\relax}|}{Aulas dialogadas e análise de filmes. Provas escritas e seminários.} \\ \hline
\multicolumn{4}{|p{\dimexpr\linewidth-2\tabcolsep-2\arrayrulewidth\relax}|}{\cellcolor{cinzaTabela}\textbf{Bibliografia Básica:}} \\ \hline
\multicolumn{4}{|p{\dimexpr\linewidth-2\tabcolsep-2\arrayrulewidth\relax}|}{1. HOBSBAWM, Eric. Era dos Extremos. São Paulo: Companhia das Letras, 2015.} \\ \hline
\multicolumn{4}{|p{\dimexpr\linewidth-2\tabcolsep-2\arrayrulewidth\relax}|}{\cellcolor{cinzaTabela}\textbf{Bibliografia Complementar:}} \\ \hline
\multicolumn{4}{|p{\dimexpr\linewidth-2\tabcolsep-2\arrayrulewidth\relax}|}{1. VAINFAS, Ronaldo. História. São Paulo: Saraiva, 2020.} \\ \hline
\end{tabularx}
\end{table}

\vspace{0.3cm}


% ---------------------------------------------------------
% TABELA EMENTA 19: GEOGRAFIA II
% ---------------------------------------------------------
\begin{table}[H]
\centering
\small
\begin{tabularx}{\linewidth}{|p{3.4cm}|X|p{3.5cm}|X|}
\hline
\multicolumn{4}{|>{\centering\arraybackslash\cellcolor{ifscgreen!20}}p{\dimexpr\linewidth-2\tabcolsep-2\arrayrulewidth\relax}|}{\textbf{\large Unidade Curricular: GEOGRAFIA II}} \\ \hline
\textbf{CH Total*:} & 80 horas & \textbf{Ano / Semestre*:} & 2º Ano (Semestres 1 e 2) \\ \hline
\textbf{CH* Prática:} & 20 horas & \textbf{CH EaD*:} & não se aplica \\ \hline
\multicolumn{2}{|l|}{\textbf{CH com Divisão de Turma*:}} & \multicolumn{2}{l|}{não se aplica} \\ \hline
\multicolumn{4}{|p{\dimexpr\linewidth-2\tabcolsep-2\arrayrulewidth\relax}|}{\cellcolor{cinzaTabela}\textbf{Objetivos:}} \\ \hline
\multicolumn{4}{|p{\dimexpr\linewidth-2\tabcolsep-2\arrayrulewidth\relax}|}{
\begin{itemize}[leftmargin=*,noitemsep,topsep=2pt]
    \item Analisar a demografia, urbanização e organização econômica do Brasil.
\end{itemize}
} \\ \hline
\multicolumn{4}{|p{\dimexpr\linewidth-2\tabcolsep-2\arrayrulewidth\relax}|}{\cellcolor{cinzaTabela}\textbf{Conteúdos:}} \\ \hline
\multicolumn{4}{|p{\dimexpr\linewidth-2\tabcolsep-2\arrayrulewidth\relax}|}{
\begin{itemize}[leftmargin=*,noitemsep,topsep=2pt]
    \item Demografia e Dinâmicas Populacionais.
    \item Urbanização e Industrialização Brasileira.
    \item Matriz Energética e Agropecuária.
\end{itemize}
} \\ \hline
\multicolumn{4}{|p{\dimexpr\linewidth-2\tabcolsep-2\arrayrulewidth\relax}|}{\cellcolor{cinzaTabela}\textbf{Metodologia de Abordagem:}} \\ \hline
\multicolumn{4}{|p{\dimexpr\linewidth-2\tabcolsep-2\arrayrulewidth\relax}|}{Aulas expositivas e análise de dados do IBGE. Avaliação por relatórios e provas.} \\ \hline
\multicolumn{4}{|p{\dimexpr\linewidth-2\tabcolsep-2\arrayrulewidth\relax}|}{\cellcolor{cinzaTabela}\textbf{Bibliografia Básica:}} \\ \hline
\multicolumn{4}{|p{\dimexpr\linewidth-2\tabcolsep-2\arrayrulewidth\relax}|}{1. SENE, Eustáquio de; MOREIRA, João Carlos. Geografia Geral e do Brasil. São Paulo: Scipione, 2018.} \\ \hline
\multicolumn{4}{|p{\dimexpr\linewidth-2\tabcolsep-2\arrayrulewidth\relax}|}{\cellcolor{cinzaTabela}\textbf{Bibliografia Complementar:}} \\ \hline
\multicolumn{4}{|p{\dimexpr\linewidth-2\tabcolsep-2\arrayrulewidth\relax}|}{1. ROSS, Jurandyr L. Sanches. Geografia do Brasil. São Paulo: EDUSP, 2017.} \\ \hline
\end{tabularx}
\end{table}

\vspace{0.3cm}


% ---------------------------------------------------------
% TABELA EMENTA 20: BIOLOGIA II
% ---------------------------------------------------------
\begin{table}[H]
\centering
\small
\begin{tabularx}{\linewidth}{|p{3.4cm}|X|p{3.5cm}|X|}
\hline
\multicolumn{4}{|>{\centering\arraybackslash\cellcolor{ifscgreen!20}}p{\dimexpr\linewidth-2\tabcolsep-2\arrayrulewidth\relax}|}{\textbf{\large Unidade Curricular: BIOLOGIA II}} \\ \hline
\textbf{CH Total*:} & 80 horas & \textbf{Ano / Semestre*:} & 2º Ano (Semestres 1 e 2) \\ \hline
\textbf{CH* Prática:} & 20 horas & \textbf{CH EaD*:} & não se aplica \\ \hline
\multicolumn{2}{|l|}{\textbf{CH com Divisão de Turma*:}} & \multicolumn{2}{l|}{não se aplica} \\ \hline
\multicolumn{4}{|p{\dimexpr\linewidth-2\tabcolsep-2\arrayrulewidth\relax}|}{\cellcolor{cinzaTabela}\textbf{Objetivos:}} \\ \hline
\multicolumn{4}{|p{\dimexpr\linewidth-2\tabcolsep-2\arrayrulewidth\relax}|}{
\begin{itemize}[leftmargin=*,noitemsep,topsep=2pt]
    \item Compreender os mecanismos de genética, evolução e microbiologia.
\end{itemize}
} \\ \hline
\multicolumn{4}{|p{\dimexpr\linewidth-2\tabcolsep-2\arrayrulewidth\relax}|}{\cellcolor{cinzaTabela}\textbf{Conteúdos:}} \\ \hline
\multicolumn{4}{|p{\dimexpr\linewidth-2\tabcolsep-2\arrayrulewidth\relax}|}{
\begin{itemize}[leftmargin=*,noitemsep,topsep=2pt]
    \item Genética Mendeliana e Biologia Molecular.
    \item Teorias Evolutivas.
    \item Microbiologia e Biotecnologia.
\end{itemize}
} \\ \hline
\multicolumn{4}{|p{\dimexpr\linewidth-2\tabcolsep-2\arrayrulewidth\relax}|}{\cellcolor{cinzaTabela}\textbf{Metodologia de Abordagem:}} \\ \hline
\multicolumn{4}{|p{\dimexpr\linewidth-2\tabcolsep-2\arrayrulewidth\relax}|}{Aulas práticas e teóricas. Provas e relatórios laboratoriais.} \\ \hline
\multicolumn{4}{|p{\dimexpr\linewidth-2\tabcolsep-2\arrayrulewidth\relax}|}{\cellcolor{cinzaTabela}\textbf{Bibliografia Básica:}} \\ \hline
\multicolumn{4}{|p{\dimexpr\linewidth-2\tabcolsep-2\arrayrulewidth\relax}|}{1. AMABIS; MARTHO. Biologia dos Organismos. São Paulo: Moderna, 2020.} \\ \hline
\multicolumn{4}{|p{\dimexpr\linewidth-2\tabcolsep-2\arrayrulewidth\relax}|}{\cellcolor{cinzaTabela}\textbf{Bibliografia Complementar:}} \\ \hline
\multicolumn{4}{|p{\dimexpr\linewidth-2\tabcolsep-2\arrayrulewidth\relax}|}{1. GRIFFITHS, Anthony J. F. Introdução à Genética. Rio de Janeiro: Guanabara Koogan, 2016.} \\ \hline
\end{tabularx}
\end{table}

\vspace{0.3cm}


% ---------------------------------------------------------
% TABELA EMENTA 21: FÍSICA II
% ---------------------------------------------------------
\begin{table}[H]
\centering
\small
\begin{tabularx}{\linewidth}{|p{3.4cm}|X|p{3.5cm}|X|}
\hline
\multicolumn{4}{|>{\centering\arraybackslash\cellcolor{ifscgreen!20}}p{\dimexpr\linewidth-2\tabcolsep-2\arrayrulewidth\relax}|}{\textbf{\large Unidade Curricular: FÍSICA II}} \\ \hline
\textbf{CH Total*:} & 40 horas & \textbf{Ano / Semestre*:} & 2º Ano (Semestres 1 e 2) \\ \hline
\textbf{CH* Prática:} & 10 horas & \textbf{CH EaD*:} & não se aplica \\ \hline
\multicolumn{2}{|l|}{\textbf{CH com Divisão de Turma*:}} & \multicolumn{2}{l|}{não se aplica} \\ \hline
\multicolumn{4}{|p{\dimexpr\linewidth-2\tabcolsep-2\arrayrulewidth\relax}|}{\cellcolor{cinzaTabela}\textbf{Objetivos:}} \\ \hline
\multicolumn{4}{|p{\dimexpr\linewidth-2\tabcolsep-2\arrayrulewidth\relax}|}{
\begin{itemize}[leftmargin=*,noitemsep,topsep=2pt]
    \item Compreender fenômenos térmicos e óptica geométrica.
\end{itemize}
} \\ \hline
\multicolumn{4}{|p{\dimexpr\linewidth-2\tabcolsep-2\arrayrulewidth\relax}|}{\cellcolor{cinzaTabela}\textbf{Conteúdos:}} \\ \hline
\multicolumn{4}{|p{\dimexpr\linewidth-2\tabcolsep-2\arrayrulewidth\relax}|}{
\begin{itemize}[leftmargin=*,noitemsep,topsep=2pt]
    \item Termometria, Calorimetria e Leis da Termodinâmica.
    \item Óptica Geométrica e Lentes.
\end{itemize}
} \\ \hline
\multicolumn{4}{|p{\dimexpr\linewidth-2\tabcolsep-2\arrayrulewidth\relax}|}{\cellcolor{cinzaTabela}\textbf{Metodologia de Abordagem:}} \\ \hline
\multicolumn{4}{|p{\dimexpr\linewidth-2\tabcolsep-2\arrayrulewidth\relax}|}{Experimentos e aulas expositivas. Provas teóricas e práticas.} \\ \hline
\multicolumn{4}{|p{\dimexpr\linewidth-2\tabcolsep-2\arrayrulewidth\relax}|}{\cellcolor{cinzaTabela}\textbf{Bibliografia Básica:}} \\ \hline
\multicolumn{4}{|p{\dimexpr\linewidth-2\tabcolsep-2\arrayrulewidth\relax}|}{1. HALLIDAY; RESNICK. Fundamentos de Física: Gravitação, Ondas e Termodinâmica. LTC, 2016.} \\ \hline
\multicolumn{4}{|p{\dimexpr\linewidth-2\tabcolsep-2\arrayrulewidth\relax}|}{\cellcolor{cinzaTabela}\textbf{Bibliografia Complementar:}} \\ \hline
\multicolumn{4}{|p{\dimexpr\linewidth-2\tabcolsep-2\arrayrulewidth\relax}|}{1. TIPLER; MOSCA. Física. LTC, 2014.} \\ \hline
\end{tabularx}
\end{table}

\vspace{0.3cm}


% ---------------------------------------------------------
% TABELA EMENTA 22: QUÍMICA II
% ---------------------------------------------------------
\begin{table}[H]
\centering
\small
\begin{tabularx}{\linewidth}{|p{3.4cm}|X|p{3.5cm}|X|}
\hline
\multicolumn{4}{|>{\centering\arraybackslash\cellcolor{ifscgreen!20}}p{\dimexpr\linewidth-2\tabcolsep-2\arrayrulewidth\relax}|}{\textbf{\large Unidade Curricular: QUÍMICA II}} \\ \hline
\textbf{CH Total*:} & 40 horas & \textbf{Ano / Semestre*:} & 2º Ano (Semestres 1 e 2) \\ \hline
\textbf{CH* Prática:} & 10 horas & \textbf{CH EaD*:} & não se aplica \\ \hline
\multicolumn{2}{|l|}{\textbf{CH com Divisão de Turma*:}} & \multicolumn{2}{l|}{não se aplica} \\ \hline
\multicolumn{4}{|p{\dimexpr\linewidth-2\tabcolsep-2\arrayrulewidth\relax}|}{\cellcolor{cinzaTabela}\textbf{Objetivos:}} \\ \hline
\multicolumn{4}{|p{\dimexpr\linewidth-2\tabcolsep-2\arrayrulewidth\relax}|}{
\begin{itemize}[leftmargin=*,noitemsep,topsep=2pt]
    \item Compreender soluções, estequiometria e termoquímica.
\end{itemize}
} \\ \hline
\multicolumn{4}{|p{\dimexpr\linewidth-2\tabcolsep-2\arrayrulewidth\relax}|}{\cellcolor{cinzaTabela}\textbf{Conteúdos:}} \\ \hline
\multicolumn{4}{|p{\dimexpr\linewidth-2\tabcolsep-2\arrayrulewidth\relax}|}{
\begin{itemize}[leftmargin=*,noitemsep,topsep=2pt]
    \item Cálculos Estequiométricos e Soluções.
    \item Termoquímica e Cinética Química.
\end{itemize}
} \\ \hline
\multicolumn{4}{|p{\dimexpr\linewidth-2\tabcolsep-2\arrayrulewidth\relax}|}{\cellcolor{cinzaTabela}\textbf{Metodologia de Abordagem:}} \\ \hline
\multicolumn{4}{|p{\dimexpr\linewidth-2\tabcolsep-2\arrayrulewidth\relax}|}{Aulas em laboratório e exercícios teóricos. Provas e relatórios.} \\ \hline
\multicolumn{4}{|p{\dimexpr\linewidth-2\tabcolsep-2\arrayrulewidth\relax}|}{\cellcolor{cinzaTabela}\textbf{Bibliografia Básica:}} \\ \hline
\multicolumn{4}{|p{\dimexpr\linewidth-2\tabcolsep-2\arrayrulewidth\relax}|}{1. RUSSELL, John B. Química Geral. Makron Books, 2018.} \\ \hline
\multicolumn{4}{|p{\dimexpr\linewidth-2\tabcolsep-2\arrayrulewidth\relax}|}{\cellcolor{cinzaTabela}\textbf{Bibliografia Complementar:}} \\ \hline
\multicolumn{4}{|p{\dimexpr\linewidth-2\tabcolsep-2\arrayrulewidth\relax}|}{1. ATKINS; JONES. Princípios de Química. Bookman, 2012.} \\ \hline
\end{tabularx}
\end{table}

\vspace{0.3cm}


% ---------------------------------------------------------
% TABELA EMENTA 23: SOCIOLOGIA I
% ---------------------------------------------------------
\begin{table}[H]
\centering
\small
\begin{tabularx}{\linewidth}{|p{3.4cm}|X|p{3.5cm}|X|}
\hline
\multicolumn{4}{|>{\centering\arraybackslash\cellcolor{ifscgreen!20}}p{\dimexpr\linewidth-2\tabcolsep-2\arrayrulewidth\relax}|}{\textbf{\large Unidade Curricular: SOCIOLOGIA I}} \\ \hline
\textbf{CH Total*:} & 80 horas & \textbf{Ano / Semestre*:} & 2º Ano (Semestres 1 e 2) \\ \hline
\textbf{CH* Prática:} & 0 horas & \textbf{CH EaD*:} & não se aplica \\ \hline
\multicolumn{2}{|l|}{\textbf{CH com Divisão de Turma*:}} & \multicolumn{2}{l|}{não se aplica} \\ \hline
\multicolumn{4}{|p{\dimexpr\linewidth-2\tabcolsep-2\arrayrulewidth\relax}|}{\cellcolor{cinzaTabela}\textbf{Objetivos:}} \\ \hline
\multicolumn{4}{|p{\dimexpr\linewidth-2\tabcolsep-2\arrayrulewidth\relax}|}{
\begin{itemize}[leftmargin=*,noitemsep,topsep=2pt]
    \item Analisar as teorias sociológicas clássicas e a formação da sociedade moderna.
\end{itemize}
} \\ \hline
\multicolumn{4}{|p{\dimexpr\linewidth-2\tabcolsep-2\arrayrulewidth\relax}|}{\cellcolor{cinzaTabela}\textbf{Conteúdos:}} \\ \hline
\multicolumn{4}{|p{\dimexpr\linewidth-2\tabcolsep-2\arrayrulewidth\relax}|}{
\begin{itemize}[leftmargin=*,noitemsep,topsep=2pt]
    \item Pensadores Clássicos: Durkheim, Marx e Weber.
    \item Estrutura Social e Desigualdade.
    \item Cultura, Identidade e Socialização.
\end{itemize}
} \\ \hline
\multicolumn{4}{|p{\dimexpr\linewidth-2\tabcolsep-2\arrayrulewidth\relax}|}{\cellcolor{cinzaTabela}\textbf{Metodologia de Abordagem:}} \\ \hline
\multicolumn{4}{|p{\dimexpr\linewidth-2\tabcolsep-2\arrayrulewidth\relax}|}{Aulas expositivo-dialogadas e leitura de artigos. Produções textuais.} \\ \hline
\multicolumn{4}{|p{\dimexpr\linewidth-2\tabcolsep-2\arrayrulewidth\relax}|}{\cellcolor{cinzaTabela}\textbf{Bibliografia Básica:}} \\ \hline
\multicolumn{4}{|p{\dimexpr\linewidth-2\tabcolsep-2\arrayrulewidth\relax}|}{1. TOMAZI, Nelson Dacio. Sociologia para o Ensino Médio. São Paulo: Saraiva, 2020.} \\ \hline
\multicolumn{4}{|p{\dimexpr\linewidth-2\tabcolsep-2\arrayrulewidth\relax}|}{\cellcolor{cinzaTabela}\textbf{Bibliografia Complementar:}} \\ \hline
\multicolumn{4}{|p{\dimexpr\linewidth-2\tabcolsep-2\arrayrulewidth\relax}|}{1. GIDDENS, Anthony. Sociologia. Porto Alegre: Penso, 2017.} \\ \hline
\end{tabularx}
\end{table}

\vspace{0.3cm}


% ---------------------------------------------------------
% TABELA EMENTA 24: ARTES
% ---------------------------------------------------------
\begin{table}[H]
\centering
\small
\begin{tabularx}{\linewidth}{|p{3.4cm}|X|p{3.5cm}|X|}
\hline
\multicolumn{4}{|>{\centering\arraybackslash\cellcolor{ifscgreen!20}}p{\dimexpr\linewidth-2\tabcolsep-2\arrayrulewidth\relax}|}{\textbf{\large Unidade Curricular: ARTES}} \\ \hline
\textbf{CH Total*:} & 80 horas & \textbf{Ano / Semestre*:} & 2º Ano (Semestres 1 e 2) \\ \hline
\textbf{CH* Prática:} & 30 horas & \textbf{CH EaD*:} & não se aplica \\ \hline
\multicolumn{2}{|l|}{\textbf{CH com Divisão de Turma*:}} & \multicolumn{2}{l|}{não se aplica} \\ \hline
\multicolumn{4}{|p{\dimexpr\linewidth-2\tabcolsep-2\arrayrulewidth\relax}|}{\cellcolor{cinzaTabela}\textbf{Objetivos:}} \\ \hline
\multicolumn{4}{|p{\dimexpr\linewidth-2\tabcolsep-2\arrayrulewidth\relax}|}{
\begin{itemize}[leftmargin=*,noitemsep,topsep=2pt]
    \item Compreender as manifestações artísticas visuais, cênicas e audiovisuais.
\end{itemize}
} \\ \hline
\multicolumn{4}{|p{\dimexpr\linewidth-2\tabcolsep-2\arrayrulewidth\relax}|}{\cellcolor{cinzaTabela}\textbf{Conteúdos:}} \\ \hline
\multicolumn{4}{|p{\dimexpr\linewidth-2\tabcolsep-2\arrayrulewidth\relax}|}{
\begin{itemize}[leftmargin=*,noitemsep,topsep=2pt]
    \item História da Arte Ocidental e Brasileira.
    \item Artes Visuais, Teatro e Dança.
    \item Arte Contemporânea e Cultura Digital.
\end{itemize}
} \\ \hline
\multicolumn{4}{|p{\dimexpr\linewidth-2\tabcolsep-2\arrayrulewidth\relax}|}{\cellcolor{cinzaTabela}\textbf{Metodologia de Abordagem:}} \\ \hline
\multicolumn{4}{|p{\dimexpr\linewidth-2\tabcolsep-2\arrayrulewidth\relax}|}{Oficinas práticas de criação artística e apreciação estética. Trabalhos em equipe.} \\ \hline
\multicolumn{4}{|p{\dimexpr\linewidth-2\tabcolsep-2\arrayrulewidth\relax}|}{\cellcolor{cinzaTabela}\textbf{Bibliografia Básica:}} \\ \hline
\multicolumn{4}{|p{\dimexpr\linewidth-2\tabcolsep-2\arrayrulewidth\relax}|}{1. STRICKLAND, Carol. Arte Comentada. Rio de Janeiro: Ediouro, 2014.} \\ \hline
\multicolumn{4}{|p{\dimexpr\linewidth-2\tabcolsep-2\arrayrulewidth\relax}|}{\cellcolor{cinzaTabela}\textbf{Bibliografia Complementar:}} \\ \hline
\multicolumn{4}{|p{\dimexpr\linewidth-2\tabcolsep-2\arrayrulewidth\relax}|}{1. GOMBRICH, E. H. A História da Arte. Rio de Janeiro: LTC, 2015.} \\ \hline
\end{tabularx}
\end{table}

\vspace{0.3cm}


% ---------------------------------------------------------
% TABELA EMENTA 25: EDUCAÇÃO FÍSICA II
% ---------------------------------------------------------
\begin{table}[H]
\centering
\small
\begin{tabularx}{\linewidth}{|p{3.4cm}|X|p{3.5cm}|X|}
\hline
\multicolumn{4}{|>{\centering\arraybackslash\cellcolor{ifscgreen!20}}p{\dimexpr\linewidth-2\tabcolsep-2\arrayrulewidth\relax}|}{\textbf{\large Unidade Curricular: EDUCAÇÃO FÍSICA II}} \\ \hline
\textbf{CH Total*:} & 80 horas & \textbf{Ano / Semestre*:} & 2º Ano (Semestres 1 e 2) \\ \hline
\textbf{CH* Prática:} & 60 horas & \textbf{CH EaD*:} & não se aplica \\ \hline
\multicolumn{2}{|l|}{\textbf{CH com Divisão de Turma*:}} & \multicolumn{2}{l|}{não se aplica} \\ \hline
\multicolumn{4}{|p{\dimexpr\linewidth-2\tabcolsep-2\arrayrulewidth\relax}|}{\cellcolor{cinzaTabela}\textbf{Objetivos:}} \\ \hline
\multicolumn{4}{|p{\dimexpr\linewidth-2\tabcolsep-2\arrayrulewidth\relax}|}{
\begin{itemize}[leftmargin=*,noitemsep,topsep=2pt]
    \item Aprofundar vivências corporais, expressão motora e esportes alternativos.
\end{itemize}
} \\ \hline
\multicolumn{4}{|p{\dimexpr\linewidth-2\tabcolsep-2\arrayrulewidth\relax}|}{\cellcolor{cinzaTabela}\textbf{Conteúdos:}} \\ \hline
\multicolumn{4}{|p{\dimexpr\linewidth-2\tabcolsep-2\arrayrulewidth\relax}|}{
\begin{itemize}[leftmargin=*,noitemsep,topsep=2pt]
    \item Ginástica, Lutas e Esportes da Natureza.
    \item Corpo, Mídia e Sociedade.
    \item Atividade Física e Qualidade de Vida.
\end{itemize}
} \\ \hline
\multicolumn{4}{|p{\dimexpr\linewidth-2\tabcolsep-2\arrayrulewidth\relax}|}{\cellcolor{cinzaTabela}\textbf{Metodologia de Abordagem:}} \\ \hline
\multicolumn{4}{|p{\dimexpr\linewidth-2\tabcolsep-2\arrayrulewidth\relax}|}{Aulas práticas no ginásio e espaços abertos. Avaliação por participação ativa.} \\ \hline
\multicolumn{4}{|p{\dimexpr\linewidth-2\tabcolsep-2\arrayrulewidth\relax}|}{\cellcolor{cinzaTabela}\textbf{Bibliografia Básica:}} \\ \hline
\multicolumn{4}{|p{\dimexpr\linewidth-2\tabcolsep-2\arrayrulewidth\relax}|}{1. DARIDO, Suraya C. Educação Física na Escola. Guanabara Koogan, 2017.} \\ \hline
\multicolumn{4}{|p{\dimexpr\linewidth-2\tabcolsep-2\arrayrulewidth\relax}|}{\cellcolor{cinzaTabela}\textbf{Bibliografia Complementar:}} \\ \hline
\multicolumn{4}{|p{\dimexpr\linewidth-2\tabcolsep-2\arrayrulewidth\relax}|}{1. KUNZ, Elenor. Transformação Didático-Pedagógica do Esporte. Unijuí, 2014.} \\ \hline
\end{tabularx}
\end{table}

\vspace{0.3cm}


% ---------------------------------------------------------
% TABELA EMENTA 26: MATEMÁTICA PARA ADMINISTRAÇÃO
% ---------------------------------------------------------
\begin{table}[H]
\centering
\small
\begin{tabularx}{\linewidth}{|p{3.4cm}|X|p{3.5cm}|X|}
\hline
\multicolumn{4}{|>{\centering\arraybackslash\cellcolor{ifscgreen!20}}p{\dimexpr\linewidth-2\tabcolsep-2\arrayrulewidth\relax}|}{\textbf{\large Unidade Curricular: MATEMÁTICA PARA ADMINISTRAÇÃO}} \\ \hline
\textbf{CH Total*:} & 80 horas & \textbf{Ano / Semestre*:} & 2º Ano (Semestre 1) \\ \hline
\textbf{CH* Prática:} & 30 horas & \textbf{CH EaD*:} & não se aplica \\ \hline
\multicolumn{2}{|l|}{\textbf{CH com Divisão de Turma*:}} & \multicolumn{2}{l|}{não se aplica} \\ \hline
\multicolumn{4}{|p{\dimexpr\linewidth-2\tabcolsep-2\arrayrulewidth\relax}|}{\cellcolor{cinzaTabela}\textbf{Objetivos:}} \\ \hline
\multicolumn{4}{|p{\dimexpr\linewidth-2\tabcolsep-2\arrayrulewidth\relax}|}{
\begin{itemize}[leftmargin=*,noitemsep,topsep=2pt]
    \item Aplicar conceitos matemáticos e estatísticos na análise de dados financeiros e administrativos.
\end{itemize}
} \\ \hline
\multicolumn{4}{|p{\dimexpr\linewidth-2\tabcolsep-2\arrayrulewidth\relax}|}{\cellcolor{cinzaTabela}\textbf{Conteúdos:}} \\ \hline
\multicolumn{4}{|p{\dimexpr\linewidth-2\tabcolsep-2\arrayrulewidth\relax}|}{
\begin{itemize}[leftmargin=*,noitemsep,topsep=2pt]
    \item Matemática Financeira Avançada.
    \item Estatística Descritiva para Gestão.
    \item Análise de Gráficos e Indicadores.
\end{itemize}
} \\ \hline
\multicolumn{4}{|p{\dimexpr\linewidth-2\tabcolsep-2\arrayrulewidth\relax}|}{\cellcolor{cinzaTabela}\textbf{Metodologia de Abordagem:}} \\ \hline
\multicolumn{4}{|p{\dimexpr\linewidth-2\tabcolsep-2\arrayrulewidth\relax}|}{Resolução de casos práticos em planilhas eletrônicas. Avaliações práticas e escritas.} \\ \hline
\multicolumn{4}{|p{\dimexpr\linewidth-2\tabcolsep-2\arrayrulewidth\relax}|}{\cellcolor{cinzaTabela}\textbf{Bibliografia Básica:}} \\ \hline
\multicolumn{4}{|p{\dimexpr\linewidth-2\tabcolsep-2\arrayrulewidth\relax}|}{1. ASSAF NETO, Alexandre. Matemática Financeira e suas Aplicações. São Paulo: Atlas, 2019.} \\ \hline
\multicolumn{4}{|p{\dimexpr\linewidth-2\tabcolsep-2\arrayrulewidth\relax}|}{\cellcolor{cinzaTabela}\textbf{Bibliografia Complementar:}} \\ \hline
\multicolumn{4}{|p{\dimexpr\linewidth-2\tabcolsep-2\arrayrulewidth\relax}|}{1. BUSSAB, Wilton O.; MORETTIN, Pedro A. Estatística Básica. São Paulo: Saraiva, 2017.} \\ \hline
\end{tabularx}
\end{table}

\vspace{0.3cm}


% ---------------------------------------------------------
% TABELA EMENTA 27: GESTÃO DE PESSOAS E RELAÇÕES NO TRABALHO
% ---------------------------------------------------------
\begin{table}[H]
\centering
\small
\begin{tabularx}{\linewidth}{|p{3.4cm}|X|p{3.5cm}|X|}
\hline
\multicolumn{4}{|>{\centering\arraybackslash\cellcolor{ifscgreen!20}}p{\dimexpr\linewidth-2\tabcolsep-2\arrayrulewidth\relax}|}{\textbf{\large Unidade Curricular: GESTÃO DE PESSOAS E RELAÇÕES NO TRABALHO}} \\ \hline
\textbf{CH Total*:} & 80 horas & \textbf{Ano / Semestre*:} & 2º Ano (Semestre 1) \\ \hline
\textbf{CH* Prática:} & 30 horas & \textbf{CH EaD*:} & não se aplica \\ \hline
\multicolumn{2}{|l|}{\textbf{CH com Divisão de Turma*:}} & \multicolumn{2}{l|}{não se aplica} \\ \hline
\multicolumn{4}{|p{\dimexpr\linewidth-2\tabcolsep-2\arrayrulewidth\relax}|}{\cellcolor{cinzaTabela}\textbf{Objetivos:}} \\ \hline
\multicolumn{4}{|p{\dimexpr\linewidth-2\tabcolsep-2\arrayrulewidth\relax}|}{
\begin{itemize}[leftmargin=*,noitemsep,topsep=2pt]
    \item Executar processos operacionais de atração, seleção e desenvolvimento de talentos;
    \item Compreender a legislação trabalhista e rotinas de folha de pagamento;
\end{itemize}
} \\ \hline
\multicolumn{4}{|p{\dimexpr\linewidth-2\tabcolsep-2\arrayrulewidth\relax}|}{\cellcolor{cinzaTabela}\textbf{Conteúdos:}} \\ \hline
\multicolumn{4}{|p{\dimexpr\linewidth-2\tabcolsep-2\arrayrulewidth\relax}|}{
\begin{itemize}[leftmargin=*,noitemsep,topsep=2pt]
    \item Subsistemas de Gestão de Pessoas (Prover, Desenvolver, Manter, Monitorar).
    \item Recrutamento, Seleção, Treinamento e Desenvolvimento.
    \item Legislação Trabalhista (CLT), Saúde e Segurança no Trabalho.
\end{itemize}
} \\ \hline
\multicolumn{4}{|p{\dimexpr\linewidth-2\tabcolsep-2\arrayrulewidth\relax}|}{\cellcolor{cinzaTabela}\textbf{Metodologia de Abordagem:}} \\ \hline
\multicolumn{4}{|p{\dimexpr\linewidth-2\tabcolsep-2\arrayrulewidth\relax}|}{Dinâmicas de grupo e simulações corporativas. Relatório de simulação de recrutamento e avaliação teórica.} \\ \hline
\multicolumn{4}{|p{\dimexpr\linewidth-2\tabcolsep-2\arrayrulewidth\relax}|}{\cellcolor{cinzaTabela}\textbf{Bibliografia Básica:}} \\ \hline
\multicolumn{4}{|p{\dimexpr\linewidth-2\tabcolsep-2\arrayrulewidth\relax}|}{1. CHIAVENATO, Idalberto. Gestão de Pessoas. 5. ed. Rio de Janeiro: Atlas, 2020.} \\ \hline
\multicolumn{4}{|p{\dimexpr\linewidth-2\tabcolsep-2\arrayrulewidth\relax}|}{\cellcolor{cinzaTabela}\textbf{Bibliografia Complementar:}} \\ \hline
\multicolumn{4}{|p{\dimexpr\linewidth-2\tabcolsep-2\arrayrulewidth\relax}|}{1. MARRAS, Jean Pierre. Gestão de Pessoas em Época de Mudanças. São Paulo: Saraiva, 2016.} \\ \hline
\end{tabularx}
\end{table}

\vspace{0.3cm}


% ---------------------------------------------------------
% TABELA EMENTA 28: GESTÃO DE MARKETING I
% ---------------------------------------------------------
\begin{table}[H]
\centering
\small
\begin{tabularx}{\linewidth}{|p{3.4cm}|X|p{3.5cm}|X|}
\hline
\multicolumn{4}{|>{\centering\arraybackslash\cellcolor{ifscgreen!20}}p{\dimexpr\linewidth-2\tabcolsep-2\arrayrulewidth\relax}|}{\textbf{\large Unidade Curricular: GESTÃO DE MARKETING I}} \\ \hline
\textbf{CH Total*:} & 80 horas & \textbf{Ano / Semestre*:} & 2º Ano (Semestre 1) \\ \hline
\textbf{CH* Prática:} & 30 horas & \textbf{CH EaD*:} & não se aplica \\ \hline
\multicolumn{2}{|l|}{\textbf{CH com Divisão de Turma*:}} & \multicolumn{2}{l|}{não se aplica} \\ \hline
\multicolumn{4}{|p{\dimexpr\linewidth-2\tabcolsep-2\arrayrulewidth\relax}|}{\cellcolor{cinzaTabela}\textbf{Objetivos:}} \\ \hline
\multicolumn{4}{|p{\dimexpr\linewidth-2\tabcolsep-2\arrayrulewidth\relax}|}{
\begin{itemize}[leftmargin=*,noitemsep,topsep=2pt]
    \item Compreender os conceitos centrais de valor, mercado e comportamento do cliente;
    \item Elaborar diagnósticos mercadológicos para empresas regionais;
\end{itemize}
} \\ \hline
\multicolumn{4}{|p{\dimexpr\linewidth-2\tabcolsep-2\arrayrulewidth\relax}|}{\cellcolor{cinzaTabela}\textbf{Conteúdos:}} \\ \hline
\multicolumn{4}{|p{\dimexpr\linewidth-2\tabcolsep-2\arrayrulewidth\relax}|}{
\begin{itemize}[leftmargin=*,noitemsep,topsep=2pt]
    \item Fundamentos de Marketing e Ambiente de Negócios.
    \item Comportamento do Consumidor e Pesquisa de Mercado.
    \item Mix de Marketing (4 Ps: Produto, Preço, Praça, Promoção).
\end{itemize}
} \\ \hline
\multicolumn{4}{|p{\dimexpr\linewidth-2\tabcolsep-2\arrayrulewidth\relax}|}{\cellcolor{cinzaTabela}\textbf{Metodologia de Abordagem:}} \\ \hline
\multicolumn{4}{|p{\dimexpr\linewidth-2\tabcolsep-2\arrayrulewidth\relax}|}{Análise de casos mercadológicos reais e desenvolvimento de pesquisas de campo. Provas e trabalho prático.} \\ \hline
\multicolumn{4}{|p{\dimexpr\linewidth-2\tabcolsep-2\arrayrulewidth\relax}|}{\cellcolor{cinzaTabela}\textbf{Bibliografia Básica:}} \\ \hline
\multicolumn{4}{|p{\dimexpr\linewidth-2\tabcolsep-2\arrayrulewidth\relax}|}{1. KOTLER, Philip; ARMSTRONG, Gary. Princípios de Marketing. 15. ed. Pearson, 2015.} \\ \hline
\multicolumn{4}{|p{\dimexpr\linewidth-2\tabcolsep-2\arrayrulewidth\relax}|}{\cellcolor{cinzaTabela}\textbf{Bibliografia Complementar:}} \\ \hline
\multicolumn{4}{|p{\dimexpr\linewidth-2\tabcolsep-2\arrayrulewidth\relax}|}{1. LAS CASAS, Alexandre Luzzi. Administração de Marketing. São Paulo: Atlas, 2019.} \\ \hline
\end{tabularx}
\end{table}

\vspace{0.3cm}


% ---------------------------------------------------------
% TABELA EMENTA 29: GESTÃO FINANCEIRA
% ---------------------------------------------------------
\begin{table}[H]
\centering
\small
\begin{tabularx}{\linewidth}{|p{3.4cm}|X|p{3.5cm}|X|}
\hline
\multicolumn{4}{|>{\centering\arraybackslash\cellcolor{ifscgreen!20}}p{\dimexpr\linewidth-2\tabcolsep-2\arrayrulewidth\relax}|}{\textbf{\large Unidade Curricular: GESTÃO FINANCEIRA}} \\ \hline
\textbf{CH Total*:} & 80 horas & \textbf{Ano / Semestre*:} & 2º Ano (Semestre 2) \\ \hline
\textbf{CH* Prática:} & 30 horas & \textbf{CH EaD*:} & não se aplica \\ \hline
\multicolumn{2}{|l|}{\textbf{CH com Divisão de Turma*:}} & \multicolumn{2}{l|}{não se aplica} \\ \hline
\multicolumn{4}{|p{\dimexpr\linewidth-2\tabcolsep-2\arrayrulewidth\relax}|}{\cellcolor{cinzaTabela}\textbf{Objetivos:}} \\ \hline
\multicolumn{4}{|p{\dimexpr\linewidth-2\tabcolsep-2\arrayrulewidth\relax}|}{
\begin{itemize}[leftmargin=*,noitemsep,topsep=2pt]
    \item Analisar demonstrativos financeiros e planejar o fluxo de caixa corporativo.
\end{itemize}
} \\ \hline
\multicolumn{4}{|p{\dimexpr\linewidth-2\tabcolsep-2\arrayrulewidth\relax}|}{\cellcolor{cinzaTabela}\textbf{Conteúdos:}} \\ \hline
\multicolumn{4}{|p{\dimexpr\linewidth-2\tabcolsep-2\arrayrulewidth\relax}|}{
\begin{itemize}[leftmargin=*,noitemsep,topsep=2pt]
    \item Demonstrações Financeiras (DRE, Balanço Patrimonial).
    \item Gestão do Fluxo de Caixa e Capital de Giro.
    \item Análise de Custos e Ponto de Equilíbrio.
\end{itemize}
} \\ \hline
\multicolumn{4}{|p{\dimexpr\linewidth-2\tabcolsep-2\arrayrulewidth\relax}|}{\cellcolor{cinzaTabela}\textbf{Metodologia de Abordagem:}} \\ \hline
\multicolumn{4}{|p{\dimexpr\linewidth-2\tabcolsep-2\arrayrulewidth\relax}|}{Aulas práticas de análise de balanços e simulação de fluxo de caixa em planilha. Avaliações escritas e práticas.} \\ \hline
\multicolumn{4}{|p{\dimexpr\linewidth-2\tabcolsep-2\arrayrulewidth\relax}|}{\cellcolor{cinzaTabela}\textbf{Bibliografia Básica:}} \\ \hline
\multicolumn{4}{|p{\dimexpr\linewidth-2\tabcolsep-2\arrayrulewidth\relax}|}{1. HOJI, Masakazu. Administração Financeira e Orçamentária. 12. ed. São Paulo: Atlas, 2017.} \\ \hline
\multicolumn{4}{|p{\dimexpr\linewidth-2\tabcolsep-2\arrayrulewidth\relax}|}{\cellcolor{cinzaTabela}\textbf{Bibliografia Complementar:}} \\ \hline
\multicolumn{4}{|p{\dimexpr\linewidth-2\tabcolsep-2\arrayrulewidth\relax}|}{1. GITMAN, Lawrence J. Princípios de Administração Financeira. São Paulo: Pearson, 2018.} \\ \hline
\end{tabularx}
\end{table}

\vspace{0.3cm}


% ---------------------------------------------------------
% TABELA EMENTA 30: GESTÃO DE MARKETING II
% ---------------------------------------------------------
\begin{table}[H]
\centering
\small
\begin{tabularx}{\linewidth}{|p{3.4cm}|X|p{3.5cm}|X|}
\hline
\multicolumn{4}{|>{\centering\arraybackslash\cellcolor{ifscgreen!20}}p{\dimexpr\linewidth-2\tabcolsep-2\arrayrulewidth\relax}|}{\textbf{\large Unidade Curricular: GESTÃO DE MARKETING II}} \\ \hline
\textbf{CH Total*:} & 80 horas & \textbf{Ano / Semestre*:} & 2º Ano (Semestre 2) \\ \hline
\textbf{CH* Prática:} & 30 horas & \textbf{CH EaD*:} & não se aplica \\ \hline
\multicolumn{2}{|l|}{\textbf{CH com Divisão de Turma*:}} & \multicolumn{2}{l|}{não se aplica} \\ \hline
\multicolumn{4}{|p{\dimexpr\linewidth-2\tabcolsep-2\arrayrulewidth\relax}|}{\cellcolor{cinzaTabela}\textbf{Objetivos:}} \\ \hline
\multicolumn{4}{|p{\dimexpr\linewidth-2\tabcolsep-2\arrayrulewidth\relax}|}{
\begin{itemize}[leftmargin=*,noitemsep,topsep=2pt]
    \item Desenvolver estratégias de marketing digital, branding e redes sociais para empresas.
\end{itemize}
} \\ \hline
\multicolumn{4}{|p{\dimexpr\linewidth-2\tabcolsep-2\arrayrulewidth\relax}|}{\cellcolor{cinzaTabela}\textbf{Conteúdos:}} \\ \hline
\multicolumn{4}{|p{\dimexpr\linewidth-2\tabcolsep-2\arrayrulewidth\relax}|}{
\begin{itemize}[leftmargin=*,noitemsep,topsep=2pt]
    \item Marketing Digital e Gestão de Mídias Sociais.
    \item Branding, Identidade de Marca e Comunicação Integrada.
    \item E-commerce e Inovação em Vendas.
\end{itemize}
} \\ \hline
\multicolumn{4}{|p{\dimexpr\linewidth-2\tabcolsep-2\arrayrulewidth\relax}|}{\cellcolor{cinzaTabela}\textbf{Metodologia de Abordagem:}} \\ \hline
\multicolumn{4}{|p{\dimexpr\linewidth-2\tabcolsep-2\arrayrulewidth\relax}|}{Criação de campanhas digitais fictícias e planejamento de mídias. Trabalho prático final.} \\ \hline
\multicolumn{4}{|p{\dimexpr\linewidth-2\tabcolsep-2\arrayrulewidth\relax}|}{\cellcolor{cinzaTabela}\textbf{Bibliografia Básica:}} \\ \hline
\multicolumn{4}{|p{\dimexpr\linewidth-2\tabcolsep-2\arrayrulewidth\relax}|}{1. KOTLER, Philip; KARTAJAYA, Hermawan; SETIAWAN, Iwan. Marketing 4.0 / 5.0. Rio de Janeiro: Sextante, 2021.} \\ \hline
\multicolumn{4}{|p{\dimexpr\linewidth-2\tabcolsep-2\arrayrulewidth\relax}|}{\cellcolor{cinzaTabela}\textbf{Bibliografia Complementar:}} \\ \hline
\multicolumn{4}{|p{\dimexpr\linewidth-2\tabcolsep-2\arrayrulewidth\relax}|}{1. TORRES, Claudio. A Bíblia do Marketing Digital. São Paulo: Novatec, 2018.} \\ \hline
\end{tabularx}
\end{table}

\vspace{0.3cm}


% ---------------------------------------------------------
% TABELA EMENTA 31: EMPREENDEDORISMO I
% ---------------------------------------------------------
\begin{table}[H]
\centering
\small
\begin{tabularx}{\linewidth}{|p{3.4cm}|X|p{3.5cm}|X|}
\hline
\multicolumn{4}{|>{\centering\arraybackslash\cellcolor{ifscgreen!20}}p{\dimexpr\linewidth-2\tabcolsep-2\arrayrulewidth\relax}|}{\textbf{\large Unidade Curricular: EMPREENDEDORISMO I}} \\ \hline
\textbf{CH Total*:} & 80 horas & \textbf{Ano / Semestre*:} & 2º Ano (Semestre 2) \\ \hline
\textbf{CH* Prática:} & 30 horas & \textbf{CH EaD*:} & não se aplica \\ \hline
\multicolumn{2}{|l|}{\textbf{CH com Divisão de Turma*:}} & \multicolumn{2}{l|}{não se aplica} \\ \hline
\multicolumn{4}{|p{\dimexpr\linewidth-2\tabcolsep-2\arrayrulewidth\relax}|}{\cellcolor{cinzaTabela}\textbf{Objetivos:}} \\ \hline
\multicolumn{4}{|p{\dimexpr\linewidth-2\tabcolsep-2\arrayrulewidth\relax}|}{
\begin{itemize}[leftmargin=*,noitemsep,topsep=2pt]
    \item Identificar oportunidades de negócios inovadores e utilizar ferramentas de modelagem (Business Model Canvas).
\end{itemize}
} \\ \hline
\multicolumn{4}{|p{\dimexpr\linewidth-2\tabcolsep-2\arrayrulewidth\relax}|}{\cellcolor{cinzaTabela}\textbf{Conteúdos:}} \\ \hline
\multicolumn{4}{|p{\dimexpr\linewidth-2\tabcolsep-2\arrayrulewidth\relax}|}{
\begin{itemize}[leftmargin=*,noitemsep,topsep=2pt]
    \item Perfil Empreendedor e Identificação de Oportunidades.
    \item Modelagem de Negócios com Canvas.
    \item Validação de Propostas de Valor e Lean Startup.
\end{itemize}
} \\ \hline
\multicolumn{4}{|p{\dimexpr\linewidth-2\tabcolsep-2\arrayrulewidth\relax}|}{\cellcolor{cinzaTabela}\textbf{Metodologia de Abordagem:}} \\ \hline
\multicolumn{4}{|p{\dimexpr\linewidth-2\tabcolsep-2\arrayrulewidth\relax}|}{Oficinas de ideação, criação de protótipos de negócios e apresentações no formato Pitch.} \\ \hline
\multicolumn{4}{|p{\dimexpr\linewidth-2\tabcolsep-2\arrayrulewidth\relax}|}{\cellcolor{cinzaTabela}\textbf{Bibliografia Básica:}} \\ \hline
\multicolumn{4}{|p{\dimexpr\linewidth-2\tabcolsep-2\arrayrulewidth\relax}|}{1. OSTERWALDER, Alexander; PIGNEUR, Yves. Business Model Generation: Inovação em Modelos de Negócios. Alta Books, 2011.} \\ \hline
\multicolumn{4}{|p{\dimexpr\linewidth-2\tabcolsep-2\arrayrulewidth\relax}|}{\cellcolor{cinzaTabela}\textbf{Bibliografia Complementar:}} \\ \hline
\multicolumn{4}{|p{\dimexpr\linewidth-2\tabcolsep-2\arrayrulewidth\relax}|}{1. DORNELAS, José. Empreendedorismo: Transformando Ideias em Negócios. Atlas, 2018.} \\ \hline
\end{tabularx}
\end{table}

\vspace{0.3cm}


% ---------------------------------------------------------
% TABELA EMENTA 32: OFICINA DE INTEGRAÇÃO II
% ---------------------------------------------------------
\begin{table}[H]
\centering
\small
\begin{tabularx}{\linewidth}{|p{3.4cm}|X|p{3.5cm}|X|}
\hline
\multicolumn{4}{|>{\centering\arraybackslash\cellcolor{ifscgreen!20}}p{\dimexpr\linewidth-2\tabcolsep-2\arrayrulewidth\relax}|}{\textbf{\large Unidade Curricular: OFICINA DE INTEGRAÇÃO II}} \\ \hline
\textbf{CH Total*:} & 40 horas & \textbf{Ano / Semestre*:} & 2º Ano (Semestres 1 e 2) \\ \hline
\textbf{CH* Prática:} & 30 horas & \textbf{CH EaD*:} & não se aplica \\ \hline
\multicolumn{2}{|l|}{\textbf{CH com Divisão de Turma*:}} & \multicolumn{2}{l|}{não se aplica} \\ \hline
\multicolumn{4}{|p{\dimexpr\linewidth-2\tabcolsep-2\arrayrulewidth\relax}|}{\cellcolor{cinzaTabela}\textbf{Objetivos:}} \\ \hline
\multicolumn{4}{|p{\dimexpr\linewidth-2\tabcolsep-2\arrayrulewidth\relax}|}{
\begin{itemize}[leftmargin=*,noitemsep,topsep=2pt]
    \item Articular saberes das disciplinas técnicas da 2ª fase em diagnósticos regionais.
\end{itemize}
} \\ \hline
\multicolumn{4}{|p{\dimexpr\linewidth-2\tabcolsep-2\arrayrulewidth\relax}|}{\cellcolor{cinzaTabela}\textbf{Conteúdos:}} \\ \hline
\multicolumn{4}{|p{\dimexpr\linewidth-2\tabcolsep-2\arrayrulewidth\relax}|}{
\begin{itemize}[leftmargin=*,noitemsep,topsep=2pt]
    \item Diagnóstico Organizacional Interdisciplinar.
    \item Análise Financeira e Mercadológica de Pequenos Negócios.
\end{itemize}
} \\ \hline
\multicolumn{4}{|p{\dimexpr\linewidth-2\tabcolsep-2\arrayrulewidth\relax}|}{\cellcolor{cinzaTabela}\textbf{Metodologia de Abordagem:}} \\ \hline
\multicolumn{4}{|p{\dimexpr\linewidth-2\tabcolsep-2\arrayrulewidth\relax}|}{Estudo de campo em empresas da região de Garopaba e relatórios em equipe.} \\ \hline
\multicolumn{4}{|p{\dimexpr\linewidth-2\tabcolsep-2\arrayrulewidth\relax}|}{\cellcolor{cinzaTabela}\textbf{Bibliografia Básica:}} \\ \hline
\multicolumn{4}{|p{\dimexpr\linewidth-2\tabcolsep-2\arrayrulewidth\relax}|}{1. KAUFMANN, Luiz. Projetos Integradores no Ensino Profissionalizante. São Paulo: Érica, 2018.} \\ \hline
\multicolumn{4}{|p{\dimexpr\linewidth-2\tabcolsep-2\arrayrulewidth\relax}|}{\cellcolor{cinzaTabela}\textbf{Bibliografia Complementar:}} \\ \hline
\multicolumn{4}{|p{\dimexpr\linewidth-2\tabcolsep-2\arrayrulewidth\relax}|}{1. GIL, Antonio Carlos. Métodos e Técnicas de Pesquisa Social. Atlas, 2019.} \\ \hline
\end{tabularx}
\end{table}

\vspace{0.3cm}


% ---------------------------------------------------------
% TABELA EMENTA 33: LÍNGUA PORTUGUESA E LITERATURA III
% ---------------------------------------------------------
\begin{table}[H]
\centering
\small
\begin{tabularx}{\linewidth}{|p{3.4cm}|X|p{3.5cm}|X|}
\hline
\multicolumn{4}{|>{\centering\arraybackslash\cellcolor{ifscgreen!20}}p{\dimexpr\linewidth-2\tabcolsep-2\arrayrulewidth\relax}|}{\textbf{\large Unidade Curricular: LÍNGUA PORTUGUESA E LITERATURA III}} \\ \hline
\textbf{CH Total*:} & 120 horas & \textbf{Ano / Semestre*:} & 3º Ano (Semestres 1 e 2) \\ \hline
\textbf{CH* Prática:} & 30 horas & \textbf{CH EaD*:} & não se aplica \\ \hline
\multicolumn{2}{|l|}{\textbf{CH com Divisão de Turma*:}} & \multicolumn{2}{l|}{não se aplica} \\ \hline
\multicolumn{4}{|p{\dimexpr\linewidth-2\tabcolsep-2\arrayrulewidth\relax}|}{\cellcolor{cinzaTabela}\textbf{Objetivos:}} \\ \hline
\multicolumn{4}{|p{\dimexpr\linewidth-2\tabcolsep-2\arrayrulewidth\relax}|}{
\begin{itemize}[leftmargin=*,noitemsep,topsep=2pt]
    \item Aprimorar a escrita formal acadêmica e a interpretação crítica avançada.
\end{itemize}
} \\ \hline
\multicolumn{4}{|p{\dimexpr\linewidth-2\tabcolsep-2\arrayrulewidth\relax}|}{\cellcolor{cinzaTabela}\textbf{Conteúdos:}} \\ \hline
\multicolumn{4}{|p{\dimexpr\linewidth-2\tabcolsep-2\arrayrulewidth\relax}|}{
\begin{itemize}[leftmargin=*,noitemsep,topsep=2pt]
    \item Redação Científica e Artigos Acadêmicos.
    \item Literatura Contemporânea em Língua Portuguesa.
    \item Análise Crítica de Discurso.
\end{itemize}
} \\ \hline
\multicolumn{4}{|p{\dimexpr\linewidth-2\tabcolsep-2\arrayrulewidth\relax}|}{\cellcolor{cinzaTabela}\textbf{Metodologia de Abordagem:}} \\ \hline
\multicolumn{4}{|p{\dimexpr\linewidth-2\tabcolsep-2\arrayrulewidth\relax}|}{Aulas práticas de escrita e bancas de apresentação. Artigo acadêmico final.} \\ \hline
\multicolumn{4}{|p{\dimexpr\linewidth-2\tabcolsep-2\arrayrulewidth\relax}|}{\cellcolor{cinzaTabela}\textbf{Bibliografia Básica:}} \\ \hline
\multicolumn{4}{|p{\dimexpr\linewidth-2\tabcolsep-2\arrayrulewidth\relax}|}{1. CEREJA, William R. Gramática e Literatura. Atual, 2020.} \\ \hline
\multicolumn{4}{|p{\dimexpr\linewidth-2\tabcolsep-2\arrayrulewidth\relax}|}{\cellcolor{cinzaTabela}\textbf{Bibliografia Complementar:}} \\ \hline
\multicolumn{4}{|p{\dimexpr\linewidth-2\tabcolsep-2\arrayrulewidth\relax}|}{1. MEDEIROS, João Bosco. Redação Científica. Atlas, 2019.} \\ \hline
\end{tabularx}
\end{table}

\vspace{0.3cm}


% ---------------------------------------------------------
% TABELA EMENTA 34: MATEMÁTICA III
% ---------------------------------------------------------
\begin{table}[H]
\centering
\small
\begin{tabularx}{\linewidth}{|p{3.4cm}|X|p{3.5cm}|X|}
\hline
\multicolumn{4}{|>{\centering\arraybackslash\cellcolor{ifscgreen!20}}p{\dimexpr\linewidth-2\tabcolsep-2\arrayrulewidth\relax}|}{\textbf{\large Unidade Curricular: MATEMÁTICA III}} \\ \hline
\textbf{CH Total*:} & 80 horas & \textbf{Ano / Semestre*:} & 3º Ano (Semestres 1 e 2) \\ \hline
\textbf{CH* Prática:} & 20 horas & \textbf{CH EaD*:} & não se aplica \\ \hline
\multicolumn{2}{|l|}{\textbf{CH com Divisão de Turma*:}} & \multicolumn{2}{l|}{não se aplica} \\ \hline
\multicolumn{4}{|p{\dimexpr\linewidth-2\tabcolsep-2\arrayrulewidth\relax}|}{\cellcolor{cinzaTabela}\textbf{Objetivos:}} \\ \hline
\multicolumn{4}{|p{\dimexpr\linewidth-2\tabcolsep-2\arrayrulewidth\relax}|}{
\begin{itemize}[leftmargin=*,noitemsep,topsep=2pt]
    \item Compreender geometria espacial e probabilidade.
\end{itemize}
} \\ \hline
\multicolumn{4}{|p{\dimexpr\linewidth-2\tabcolsep-2\arrayrulewidth\relax}|}{\cellcolor{cinzaTabela}\textbf{Conteúdos:}} \\ \hline
\multicolumn{4}{|p{\dimexpr\linewidth-2\tabcolsep-2\arrayrulewidth\relax}|}{
\begin{itemize}[leftmargin=*,noitemsep,topsep=2pt]
    \item Geometria Espacial de Posição e Métrica.
    \item Análise Combinatória e Probabilidade.
    \item Estatística Inferencial Básica.
\end{itemize}
} \\ \hline
\multicolumn{4}{|p{\dimexpr\linewidth-2\tabcolsep-2\arrayrulewidth\relax}|}{\cellcolor{cinzaTabela}\textbf{Metodologia de Abordagem:}} \\ \hline
\multicolumn{4}{|p{\dimexpr\linewidth-2\tabcolsep-2\arrayrulewidth\relax}|}{Aulas expositivas e resolução de problemas práticos. Provas teóricas.} \\ \hline
\multicolumn{4}{|p{\dimexpr\linewidth-2\tabcolsep-2\arrayrulewidth\relax}|}{\cellcolor{cinzaTabela}\textbf{Bibliografia Básica:}} \\ \hline
\multicolumn{4}{|p{\dimexpr\linewidth-2\tabcolsep-2\arrayrulewidth\relax}|}{1. DANTE, Luiz Roberto. Matemática. Ática, 2020.} \\ \hline
\multicolumn{4}{|p{\dimexpr\linewidth-2\tabcolsep-2\arrayrulewidth\relax}|}{\cellcolor{cinzaTabela}\textbf{Bibliografia Complementar:}} \\ \hline
\multicolumn{4}{|p{\dimexpr\linewidth-2\tabcolsep-2\arrayrulewidth\relax}|}{1. IEZZI, Gelson. Fundamentos de Matemática Elementar. Atual, 2018.} \\ \hline
\end{tabularx}
\end{table}

\vspace{0.3cm}


% ---------------------------------------------------------
% TABELA EMENTA 35: HISTÓRIA III
% ---------------------------------------------------------
\begin{table}[H]
\centering
\small
\begin{tabularx}{\linewidth}{|p{3.4cm}|X|p{3.5cm}|X|}
\hline
\multicolumn{4}{|>{\centering\arraybackslash\cellcolor{ifscgreen!20}}p{\dimexpr\linewidth-2\tabcolsep-2\arrayrulewidth\relax}|}{\textbf{\large Unidade Curricular: HISTÓRIA III}} \\ \hline
\textbf{CH Total*:} & 40 horas & \textbf{Ano / Semestre*:} & 3º Ano (Semestres 1 e 2) \\ \hline
\textbf{CH* Prática:} & 0 horas & \textbf{CH EaD*:} & não se aplica \\ \hline
\multicolumn{2}{|l|}{\textbf{CH com Divisão de Turma*:}} & \multicolumn{2}{l|}{não se aplica} \\ \hline
\multicolumn{4}{|p{\dimexpr\linewidth-2\tabcolsep-2\arrayrulewidth\relax}|}{\cellcolor{cinzaTabela}\textbf{Objetivos:}} \\ \hline
\multicolumn{4}{|p{\dimexpr\linewidth-2\tabcolsep-2\arrayrulewidth\relax}|}{
\begin{itemize}[leftmargin=*,noitemsep,topsep=2pt]
    \item Analisar a história recente do Brasil e o mundo globalizado.
\end{itemize}
} \\ \hline
\multicolumn{4}{|p{\dimexpr\linewidth-2\tabcolsep-2\arrayrulewidth\relax}|}{\cellcolor{cinzaTabela}\textbf{Conteúdos:}} \\ \hline
\multicolumn{4}{|p{\dimexpr\linewidth-2\tabcolsep-2\arrayrulewidth\relax}|}{
\begin{itemize}[leftmargin=*,noitemsep,topsep=2pt]
    \item Ditadura Militar no Brasil e Redemocratização.
    \item Globalização e Nova Ordem Mundial.
\end{itemize}
} \\ \hline
\multicolumn{4}{|p{\dimexpr\linewidth-2\tabcolsep-2\arrayrulewidth\relax}|}{\cellcolor{cinzaTabela}\textbf{Metodologia de Abordagem:}} \\ \hline
\multicolumn{4}{|p{\dimexpr\linewidth-2\tabcolsep-2\arrayrulewidth\relax}|}{Debates e produções textuais analíticas. Provas e seminários.} \\ \hline
\multicolumn{4}{|p{\dimexpr\linewidth-2\tabcolsep-2\arrayrulewidth\relax}|}{\cellcolor{cinzaTabela}\textbf{Bibliografia Básica:}} \\ \hline
\multicolumn{4}{|p{\dimexpr\linewidth-2\tabcolsep-2\arrayrulewidth\relax}|}{1. VAINFAS, Ronaldo. História Conectada. Saraiva, 2020.} \\ \hline
\multicolumn{4}{|p{\dimexpr\linewidth-2\tabcolsep-2\arrayrulewidth\relax}|}{\cellcolor{cinzaTabela}\textbf{Bibliografia Complementar:}} \\ \hline
\multicolumn{4}{|p{\dimexpr\linewidth-2\tabcolsep-2\arrayrulewidth\relax}|}{1. HOBSBAWM, Eric. Tempos Interessantes. Companhia das Letras, 2014.} \\ \hline
\end{tabularx}
\end{table}

\vspace{0.3cm}


% ---------------------------------------------------------
% TABELA EMENTA 36: GEOGRAFIA III
% ---------------------------------------------------------
\begin{table}[H]
\centering
\small
\begin{tabularx}{\linewidth}{|p{3.4cm}|X|p{3.5cm}|X|}
\hline
\multicolumn{4}{|>{\centering\arraybackslash\cellcolor{ifscgreen!20}}p{\dimexpr\linewidth-2\tabcolsep-2\arrayrulewidth\relax}|}{\textbf{\large Unidade Curricular: GEOGRAFIA III}} \\ \hline
\textbf{CH Total*:} & 40 horas & \textbf{Ano / Semestre*:} & 3º Ano (Semestres 1 e 2) \\ \hline
\textbf{CH* Prática:} & 10 horas & \textbf{CH EaD*:} & não se aplica \\ \hline
\multicolumn{2}{|l|}{\textbf{CH com Divisão de Turma*:}} & \multicolumn{2}{l|}{não se aplica} \\ \hline
\multicolumn{4}{|p{\dimexpr\linewidth-2\tabcolsep-2\arrayrulewidth\relax}|}{\cellcolor{cinzaTabela}\textbf{Objetivos:}} \\ \hline
\multicolumn{4}{|p{\dimexpr\linewidth-2\tabcolsep-2\arrayrulewidth\relax}|}{
\begin{itemize}[leftmargin=*,noitemsep,topsep=2pt]
    \item Compreender os blocos econômicos e a geografia do desenvolvimento.
\end{itemize}
} \\ \hline
\multicolumn{4}{|p{\dimexpr\linewidth-2\tabcolsep-2\arrayrulewidth\relax}|}{\cellcolor{cinzaTabela}\textbf{Conteúdos:}} \\ \hline
\multicolumn{4}{|p{\dimexpr\linewidth-2\tabcolsep-2\arrayrulewidth\relax}|}{
\begin{itemize}[leftmargin=*,noitemsep,topsep=2pt]
    \item Blocos Econômicos Mundiais.
    \item Geografia Agrária e Questões Ambientais Contemporâneas.
\end{itemize}
} \\ \hline
\multicolumn{4}{|p{\dimexpr\linewidth-2\tabcolsep-2\arrayrulewidth\relax}|}{\cellcolor{cinzaTabela}\textbf{Metodologia de Abordagem:}} \\ \hline
\multicolumn{4}{|p{\dimexpr\linewidth-2\tabcolsep-2\arrayrulewidth\relax}|}{Aulas expositivas e análise de conjuntura internacional. Provas.} \\ \hline
\multicolumn{4}{|p{\dimexpr\linewidth-2\tabcolsep-2\arrayrulewidth\relax}|}{\cellcolor{cinzaTabela}\textbf{Bibliografia Básica:}} \\ \hline
\multicolumn{4}{|p{\dimexpr\linewidth-2\tabcolsep-2\arrayrulewidth\relax}|}{1. LUCCI, Elian Alabi. Território e Sociedade. Saraiva, 2020.} \\ \hline
\multicolumn{4}{|p{\dimexpr\linewidth-2\tabcolsep-2\arrayrulewidth\relax}|}{\cellcolor{cinzaTabela}\textbf{Bibliografia Complementar:}} \\ \hline
\multicolumn{4}{|p{\dimexpr\linewidth-2\tabcolsep-2\arrayrulewidth\relax}|}{1. SANTOS, Milton. Por uma Outra Globalização. Record, 2015.} \\ \hline
\end{tabularx}
\end{table}

\vspace{0.3cm}


% ---------------------------------------------------------
% TABELA EMENTA 37: FILOSOFIA II
% ---------------------------------------------------------
\begin{table}[H]
\centering
\small
\begin{tabularx}{\linewidth}{|p{3.4cm}|X|p{3.5cm}|X|}
\hline
\multicolumn{4}{|>{\centering\arraybackslash\cellcolor{ifscgreen!20}}p{\dimexpr\linewidth-2\tabcolsep-2\arrayrulewidth\relax}|}{\textbf{\large Unidade Curricular: FILOSOFIA II}} \\ \hline
\textbf{CH Total*:} & 40 horas & \textbf{Ano / Semestre*:} & 3º Ano (Semestres 1 e 2) \\ \hline
\textbf{CH* Prática:} & 0 horas & \textbf{CH EaD*:} & não se aplica \\ \hline
\multicolumn{2}{|l|}{\textbf{CH com Divisão de Turma*:}} & \multicolumn{2}{l|}{não se aplica} \\ \hline
\multicolumn{4}{|p{\dimexpr\linewidth-2\tabcolsep-2\arrayrulewidth\relax}|}{\cellcolor{cinzaTabela}\textbf{Objetivos:}} \\ \hline
\multicolumn{4}{|p{\dimexpr\linewidth-2\tabcolsep-2\arrayrulewidth\relax}|}{
\begin{itemize}[leftmargin=*,noitemsep,topsep=2pt]
    \item Refletir sobre bioética, filosofia da ciência e estético-política.
\end{itemize}
} \\ \hline
\multicolumn{4}{|p{\dimexpr\linewidth-2\tabcolsep-2\arrayrulewidth\relax}|}{\cellcolor{cinzaTabela}\textbf{Conteúdos:}} \\ \hline
\multicolumn{4}{|p{\dimexpr\linewidth-2\tabcolsep-2\arrayrulewidth\relax}|}{
\begin{itemize}[leftmargin=*,noitemsep,topsep=2pt]
    \item Bioética e Filosofia da Tecnologia.
    \item Estética Filosófica e Sociedade do Espetáculo.
\end{itemize}
} \\ \hline
\multicolumn{4}{|p{\dimexpr\linewidth-2\tabcolsep-2\arrayrulewidth\relax}|}{\cellcolor{cinzaTabela}\textbf{Metodologia de Abordagem:}} \\ \hline
\multicolumn{4}{|p{\dimexpr\linewidth-2\tabcolsep-2\arrayrulewidth\relax}|}{Leitura de textos filosóficos e debates. Ensaio filosófico.} \\ \hline
\multicolumn{4}{|p{\dimexpr\linewidth-2\tabcolsep-2\arrayrulewidth\relax}|}{\cellcolor{cinzaTabela}\textbf{Bibliografia Básica:}} \\ \hline
\multicolumn{4}{|p{\dimexpr\linewidth-2\tabcolsep-2\arrayrulewidth\relax}|}{1. CHAUÍ, Marilena. Convite à Filosofia. Ática, 2020.} \\ \hline
\multicolumn{4}{|p{\dimexpr\linewidth-2\tabcolsep-2\arrayrulewidth\relax}|}{\cellcolor{cinzaTabela}\textbf{Bibliografia Complementar:}} \\ \hline
\multicolumn{4}{|p{\dimexpr\linewidth-2\tabcolsep-2\arrayrulewidth\relax}|}{1. ARANHA, Maria Lúcia. Filosofando. Moderna, 2019.} \\ \hline
\end{tabularx}
\end{table}

\vspace{0.3cm}


% ---------------------------------------------------------
% TABELA EMENTA 38: SOCIOLOGIA II
% ---------------------------------------------------------
\begin{table}[H]
\centering
\small
\begin{tabularx}{\linewidth}{|p{3.4cm}|X|p{3.5cm}|X|}
\hline
\multicolumn{4}{|>{\centering\arraybackslash\cellcolor{ifscgreen!20}}p{\dimexpr\linewidth-2\tabcolsep-2\arrayrulewidth\relax}|}{\textbf{\large Unidade Curricular: SOCIOLOGIA II}} \\ \hline
\textbf{CH Total*:} & 40 horas & \textbf{Ano / Semestre*:} & 3º Ano (Semestres 1 e 2) \\ \hline
\textbf{CH* Prática:} & 0 horas & \textbf{CH EaD*:} & não se aplica \\ \hline
\multicolumn{2}{|l|}{\textbf{CH com Divisão de Turma*:}} & \multicolumn{2}{l|}{não se aplica} \\ \hline
\multicolumn{4}{|p{\dimexpr\linewidth-2\tabcolsep-2\arrayrulewidth\relax}|}{\cellcolor{cinzaTabela}\textbf{Objetivos:}} \\ \hline
\multicolumn{4}{|p{\dimexpr\linewidth-2\tabcolsep-2\arrayrulewidth\relax}|}{
\begin{itemize}[leftmargin=*,noitemsep,topsep=2pt]
    \item Analisar os movimentos sociais, cidadania e sociologia brasileira.
\end{itemize}
} \\ \hline
\multicolumn{4}{|p{\dimexpr\linewidth-2\tabcolsep-2\arrayrulewidth\relax}|}{\cellcolor{cinzaTabela}\textbf{Conteúdos:}} \\ \hline
\multicolumn{4}{|p{\dimexpr\linewidth-2\tabcolsep-2\arrayrulewidth\relax}|}{
\begin{itemize}[leftmargin=*,noitemsep,topsep=2pt]
    \item Pensamento Social Brasileiro.
    \item Movimentos Sociais e Direitos de Cidadania.
\end{itemize}
} \\ \hline
\multicolumn{4}{|p{\dimexpr\linewidth-2\tabcolsep-2\arrayrulewidth\relax}|}{\cellcolor{cinzaTabela}\textbf{Metodologia de Abordagem:}} \\ \hline
\multicolumn{4}{|p{\dimexpr\linewidth-2\tabcolsep-2\arrayrulewidth\relax}|}{Aulas dialogadas e seminários em grupo.} \\ \hline
\multicolumn{4}{|p{\dimexpr\linewidth-2\tabcolsep-2\arrayrulewidth\relax}|}{\cellcolor{cinzaTabela}\textbf{Bibliografia Básica:}} \\ \hline
\multicolumn{4}{|p{\dimexpr\linewidth-2\tabcolsep-2\arrayrulewidth\relax}|}{1. TOMAZI, Nelson Dacio. Sociologia. Saraiva, 2020.} \\ \hline
\multicolumn{4}{|p{\dimexpr\linewidth-2\tabcolsep-2\arrayrulewidth\relax}|}{\cellcolor{cinzaTabela}\textbf{Bibliografia Complementar:}} \\ \hline
\multicolumn{4}{|p{\dimexpr\linewidth-2\tabcolsep-2\arrayrulewidth\relax}|}{1. FREYRE, Gilberto. Casa-Grande \& Senzala. Global, 2019.} \\ \hline
\end{tabularx}
\end{table}

\vspace{0.3cm}


% ---------------------------------------------------------
% TABELA EMENTA 39: FÍSICA III
% ---------------------------------------------------------
\begin{table}[H]
\centering
\small
\begin{tabularx}{\linewidth}{|p{3.4cm}|X|p{3.5cm}|X|}
\hline
\multicolumn{4}{|>{\centering\arraybackslash\cellcolor{ifscgreen!20}}p{\dimexpr\linewidth-2\tabcolsep-2\arrayrulewidth\relax}|}{\textbf{\large Unidade Curricular: FÍSICA III}} \\ \hline
\textbf{CH Total*:} & 80 horas & \textbf{Ano / Semestre*:} & 3º Ano (Semestres 1 e 2) \\ \hline
\textbf{CH* Prática:} & 20 horas & \textbf{CH EaD*:} & não se aplica \\ \hline
\multicolumn{2}{|l|}{\textbf{CH com Divisão de Turma*:}} & \multicolumn{2}{l|}{não se aplica} \\ \hline
\multicolumn{4}{|p{\dimexpr\linewidth-2\tabcolsep-2\arrayrulewidth\relax}|}{\cellcolor{cinzaTabela}\textbf{Objetivos:}} \\ \hline
\multicolumn{4}{|p{\dimexpr\linewidth-2\tabcolsep-2\arrayrulewidth\relax}|}{
\begin{itemize}[leftmargin=*,noitemsep,topsep=2pt]
    \item Compreender eletricidade, magnetismo e física moderna.
\end{itemize}
} \\ \hline
\multicolumn{4}{|p{\dimexpr\linewidth-2\tabcolsep-2\arrayrulewidth\relax}|}{\cellcolor{cinzaTabela}\textbf{Conteúdos:}} \\ \hline
\multicolumn{4}{|p{\dimexpr\linewidth-2\tabcolsep-2\arrayrulewidth\relax}|}{
\begin{itemize}[leftmargin=*,noitemsep,topsep=2pt]
    \item Eletrostática, Eletrodinâmica e Circuitos.
    \item Eletromagnetismo e Física Moderna Básica.
\end{itemize}
} \\ \hline
\multicolumn{4}{|p{\dimexpr\linewidth-2\tabcolsep-2\arrayrulewidth\relax}|}{\cellcolor{cinzaTabela}\textbf{Metodologia de Abordagem:}} \\ \hline
\multicolumn{4}{|p{\dimexpr\linewidth-2\tabcolsep-2\arrayrulewidth\relax}|}{Montagem de circuitos simples e aulas teóricas. Provas e práticas.} \\ \hline
\multicolumn{4}{|p{\dimexpr\linewidth-2\tabcolsep-2\arrayrulewidth\relax}|}{\cellcolor{cinzaTabela}\textbf{Bibliografia Básica:}} \\ \hline
\multicolumn{4}{|p{\dimexpr\linewidth-2\tabcolsep-2\arrayrulewidth\relax}|}{1. HALLIDAY; RESNICK. Fundamentos de Física: Eletromagnetismo. LTC, 2016.} \\ \hline
\multicolumn{4}{|p{\dimexpr\linewidth-2\tabcolsep-2\arrayrulewidth\relax}|}{\cellcolor{cinzaTabela}\textbf{Bibliografia Complementar:}} \\ \hline
\multicolumn{4}{|p{\dimexpr\linewidth-2\tabcolsep-2\arrayrulewidth\relax}|}{1. TIPLER; MOSCA. Física. LTC, 2014.} \\ \hline
\end{tabularx}
\end{table}

\vspace{0.3cm}


% ---------------------------------------------------------
% TABELA EMENTA 40: QUÍMICA III
% ---------------------------------------------------------
\begin{table}[H]
\centering
\small
\begin{tabularx}{\linewidth}{|p{3.4cm}|X|p{3.5cm}|X|}
\hline
\multicolumn{4}{|>{\centering\arraybackslash\cellcolor{ifscgreen!20}}p{\dimexpr\linewidth-2\tabcolsep-2\arrayrulewidth\relax}|}{\textbf{\large Unidade Curricular: QUÍMICA III}} \\ \hline
\textbf{CH Total*:} & 80 horas & \textbf{Ano / Semestre*:} & 3º Ano (Semestres 1 e 2) \\ \hline
\textbf{CH* Prática:} & 20 horas & \textbf{CH EaD*:} & não se aplica \\ \hline
\multicolumn{2}{|l|}{\textbf{CH com Divisão de Turma*:}} & \multicolumn{2}{l|}{não se aplica} \\ \hline
\multicolumn{4}{|p{\dimexpr\linewidth-2\tabcolsep-2\arrayrulewidth\relax}|}{\cellcolor{cinzaTabela}\textbf{Objetivos:}} \\ \hline
\multicolumn{4}{|p{\dimexpr\linewidth-2\tabcolsep-2\arrayrulewidth\relax}|}{
\begin{itemize}[leftmargin=*,noitemsep,topsep=2pt]
    \item Compreender química orgânica e eletroquímica.
\end{itemize}
} \\ \hline
\multicolumn{4}{|p{\dimexpr\linewidth-2\tabcolsep-2\arrayrulewidth\relax}|}{\cellcolor{cinzaTabela}\textbf{Conteúdos:}} \\ \hline
\multicolumn{4}{|p{\dimexpr\linewidth-2\tabcolsep-2\arrayrulewidth\relax}|}{
\begin{itemize}[leftmargin=*,noitemsep,topsep=2pt]
    \item Química Orgânica (Funções e Isomeria).
    \item Eletroquímica (Pilhas e Eletrólise).
\end{itemize}
} \\ \hline
\multicolumn{4}{|p{\dimexpr\linewidth-2\tabcolsep-2\arrayrulewidth\relax}|}{\cellcolor{cinzaTabela}\textbf{Metodologia de Abordagem:}} \\ \hline
\multicolumn{4}{|p{\dimexpr\linewidth-2\tabcolsep-2\arrayrulewidth\relax}|}{Aulas teóricas e experimentos em laboratório. Provas e relatórios.} \\ \hline
\multicolumn{4}{|p{\dimexpr\linewidth-2\tabcolsep-2\arrayrulewidth\relax}|}{\cellcolor{cinzaTabela}\textbf{Bibliografia Básica:}} \\ \hline
\multicolumn{4}{|p{\dimexpr\linewidth-2\tabcolsep-2\arrayrulewidth\relax}|}{1. RUSSELL, John B. Química Geral. Makron Books, 2018.} \\ \hline
\multicolumn{4}{|p{\dimexpr\linewidth-2\tabcolsep-2\arrayrulewidth\relax}|}{\cellcolor{cinzaTabela}\textbf{Bibliografia Complementar:}} \\ \hline
\multicolumn{4}{|p{\dimexpr\linewidth-2\tabcolsep-2\arrayrulewidth\relax}|}{1. ATKINS; JONES. Princípios de Química. Bookman, 2012.} \\ \hline
\end{tabularx}
\end{table}

\vspace{0.3cm}


% ---------------------------------------------------------
% TABELA EMENTA 41: INGLÊS
% ---------------------------------------------------------
\begin{table}[H]
\centering
\small
\begin{tabularx}{\linewidth}{|p{3.4cm}|X|p{3.5cm}|X|}
\hline
\multicolumn{4}{|>{\centering\arraybackslash\cellcolor{ifscgreen!20}}p{\dimexpr\linewidth-2\tabcolsep-2\arrayrulewidth\relax}|}{\textbf{\large Unidade Curricular: INGLÊS}} \\ \hline
\textbf{CH Total*:} & 80 horas & \textbf{Ano / Semestre*:} & 3º Ano (Semestres 1 e 2) \\ \hline
\textbf{CH* Prática:} & 20 horas & \textbf{CH EaD*:} & não se aplica \\ \hline
\multicolumn{2}{|l|}{\textbf{CH com Divisão de Turma*:}} & \multicolumn{2}{l|}{não se aplica} \\ \hline
\multicolumn{4}{|p{\dimexpr\linewidth-2\tabcolsep-2\arrayrulewidth\relax}|}{\cellcolor{cinzaTabela}\textbf{Objetivos:}} \\ \hline
\multicolumn{4}{|p{\dimexpr\linewidth-2\tabcolsep-2\arrayrulewidth\relax}|}{
\begin{itemize}[leftmargin=*,noitemsep,topsep=2pt]
    \item Desenvolver leitura instrumental e comunicação técnica em inglês corporativo.
\end{itemize}
} \\ \hline
\multicolumn{4}{|p{\dimexpr\linewidth-2\tabcolsep-2\arrayrulewidth\relax}|}{\cellcolor{cinzaTabela}\textbf{Conteúdos:}} \\ \hline
\multicolumn{4}{|p{\dimexpr\linewidth-2\tabcolsep-2\arrayrulewidth\relax}|}{
\begin{itemize}[leftmargin=*,noitemsep,topsep=2pt]
    \item Inglês Instrumental para Negócios (Business English).
    \item Vocabulário Técnico de Gestão e Tecnologia.
    \item Redação de E-mails Executivos.
\end{itemize}
} \\ \hline
\multicolumn{4}{|p{\dimexpr\linewidth-2\tabcolsep-2\arrayrulewidth\relax}|}{\cellcolor{cinzaTabela}\textbf{Metodologia de Abordagem:}} \\ \hline
\multicolumn{4}{|p{\dimexpr\linewidth-2\tabcolsep-2\arrayrulewidth\relax}|}{Leitura de artigos de negócios em inglês e exercícios práticos. Provas.} \\ \hline
\multicolumn{4}{|p{\dimexpr\linewidth-2\tabcolsep-2\arrayrulewidth\relax}|}{\cellcolor{cinzaTabela}\textbf{Bibliografia Básica:}} \\ \hline
\multicolumn{4}{|p{\dimexpr\linewidth-2\tabcolsep-2\arrayrulewidth\relax}|}{1. LATHAM-KOENIG, Christina. English File Business. Oxford, 2019.} \\ \hline
\multicolumn{4}{|p{\dimexpr\linewidth-2\tabcolsep-2\arrayrulewidth\relax}|}{\cellcolor{cinzaTabela}\textbf{Bibliografia Complementar:}} \\ \hline
\multicolumn{4}{|p{\dimexpr\linewidth-2\tabcolsep-2\arrayrulewidth\relax}|}{1. MARKS, Jonathan. Check Your English Vocabulary for Business. A\&C Black, 2017.} \\ \hline
\end{tabularx}
\end{table}

\vspace{0.3cm}


% ---------------------------------------------------------
% TABELA EMENTA 42: ARTES II
% ---------------------------------------------------------
\begin{table}[H]
\centering
\small
\begin{tabularx}{\linewidth}{|p{3.4cm}|X|p{3.5cm}|X|}
\hline
\multicolumn{4}{|>{\centering\arraybackslash\cellcolor{ifscgreen!20}}p{\dimexpr\linewidth-2\tabcolsep-2\arrayrulewidth\relax}|}{\textbf{\large Unidade Curricular: ARTES II}} \\ \hline
\textbf{CH Total*:} & 60 horas & \textbf{Ano / Semestre*:} & 3º Ano (Semestres 1 e 2) \\ \hline
\textbf{CH* Prática:} & 20 horas & \textbf{CH EaD*:} & não se aplica \\ \hline
\multicolumn{2}{|l|}{\textbf{CH com Divisão de Turma*:}} & \multicolumn{2}{l|}{não se aplica} \\ \hline
\multicolumn{4}{|p{\dimexpr\linewidth-2\tabcolsep-2\arrayrulewidth\relax}|}{\cellcolor{cinzaTabela}\textbf{Objetivos:}} \\ \hline
\multicolumn{4}{|p{\dimexpr\linewidth-2\tabcolsep-2\arrayrulewidth\relax}|}{
\begin{itemize}[leftmargin=*,noitemsep,topsep=2pt]
    \item Explorar o design, patrimônio cultural e indústrias criativas.
\end{itemize}
} \\ \hline
\multicolumn{4}{|p{\dimexpr\linewidth-2\tabcolsep-2\arrayrulewidth\relax}|}{\cellcolor{cinzaTabela}\textbf{Conteúdos:}} \\ \hline
\multicolumn{4}{|p{\dimexpr\linewidth-2\tabcolsep-2\arrayrulewidth\relax}|}{
\begin{itemize}[leftmargin=*,noitemsep,topsep=2pt]
    \item Patrimônio Cultural e Museologia.
    \item Design, Fotografia e Audiovisual.
    \item Indústrias Criativas.
\end{itemize}
} \\ \hline
\multicolumn{4}{|p{\dimexpr\linewidth-2\tabcolsep-2\arrayrulewidth\relax}|}{\cellcolor{cinzaTabela}\textbf{Metodologia de Abordagem:}} \\ \hline
\multicolumn{4}{|p{\dimexpr\linewidth-2\tabcolsep-2\arrayrulewidth\relax}|}{Projetos de produção audiovisual e exposições. Avaliação continuada.} \\ \hline
\multicolumn{4}{|p{\dimexpr\linewidth-2\tabcolsep-2\arrayrulewidth\relax}|}{\cellcolor{cinzaTabela}\textbf{Bibliografia Básica:}} \\ \hline
\multicolumn{4}{|p{\dimexpr\linewidth-2\tabcolsep-2\arrayrulewidth\relax}|}{1. STRICKLAND, Carol. Arte Comentada. Ediouro, 2014.} \\ \hline
\multicolumn{4}{|p{\dimexpr\linewidth-2\tabcolsep-2\arrayrulewidth\relax}|}{\cellcolor{cinzaTabela}\textbf{Bibliografia Complementar:}} \\ \hline
\multicolumn{4}{|p{\dimexpr\linewidth-2\tabcolsep-2\arrayrulewidth\relax}|}{1. GOMBRICH, E. H. A História da Arte. LTC, 2015.} \\ \hline
\end{tabularx}
\end{table}

\vspace{0.3cm}


% ---------------------------------------------------------
% TABELA EMENTA 43: GESTÃO DE OPERAÇÕES E QUALIDADE
% ---------------------------------------------------------
\begin{table}[H]
\centering
\small
\begin{tabularx}{\linewidth}{|p{3.4cm}|X|p{3.5cm}|X|}
\hline
\multicolumn{4}{|>{\centering\arraybackslash\cellcolor{ifscgreen!20}}p{\dimexpr\linewidth-2\tabcolsep-2\arrayrulewidth\relax}|}{\textbf{\large Unidade Curricular: GESTÃO DE OPERAÇÕES E QUALIDADE}} \\ \hline
\textbf{CH Total*:} & 80 horas & \textbf{Ano / Semestre*:} & 3º Ano (Semestre 1) \\ \hline
\textbf{CH* Prática:} & 30 horas & \textbf{CH EaD*:} & não se aplica \\ \hline
\multicolumn{2}{|l|}{\textbf{CH com Divisão de Turma*:}} & \multicolumn{2}{l|}{não se aplica} \\ \hline
\multicolumn{4}{|p{\dimexpr\linewidth-2\tabcolsep-2\arrayrulewidth\relax}|}{\cellcolor{cinzaTabela}\textbf{Objetivos:}} \\ \hline
\multicolumn{4}{|p{\dimexpr\linewidth-2\tabcolsep-2\arrayrulewidth\relax}|}{
\begin{itemize}[leftmargin=*,noitemsep,topsep=2pt]
    \item Compreender os processos produtivos e aplicar ferramentas de gestão da qualidade nas empresas.
\end{itemize}
} \\ \hline
\multicolumn{4}{|p{\dimexpr\linewidth-2\tabcolsep-2\arrayrulewidth\relax}|}{\cellcolor{cinzaTabela}\textbf{Conteúdos:}} \\ \hline
\multicolumn{4}{|p{\dimexpr\linewidth-2\tabcolsep-2\arrayrulewidth\relax}|}{
\begin{itemize}[leftmargin=*,noitemsep,topsep=2pt]
    \item Introdução à Gestão de Operações e Serviços.
    \item Planejamento de Estoques e Logística Operacional.
    \item Ferramentas Clássicas da Qualidade (PDCA, Ishikawa, 5S, Kaizen).
\end{itemize}
} \\ \hline
\multicolumn{4}{|p{\dimexpr\linewidth-2\tabcolsep-2\arrayrulewidth\relax}|}{\cellcolor{cinzaTabela}\textbf{Metodologia de Abordagem:}} \\ \hline
\multicolumn{4}{|p{\dimexpr\linewidth-2\tabcolsep-2\arrayrulewidth\relax}|}{Aulas práticas de aplicação de ferramentas da qualidade em casos reais. Projeto 5S/Ishikawa.} \\ \hline
\multicolumn{4}{|p{\dimexpr\linewidth-2\tabcolsep-2\arrayrulewidth\relax}|}{\cellcolor{cinzaTabela}\textbf{Bibliografia Básica:}} \\ \hline
\multicolumn{4}{|p{\dimexpr\linewidth-2\tabcolsep-2\arrayrulewidth\relax}|}{1. SLACK, Nigel et al. Administração da Produção. 8. ed. São Paulo: Atlas, 2018.} \\ \hline
\multicolumn{4}{|p{\dimexpr\linewidth-2\tabcolsep-2\arrayrulewidth\relax}|}{\cellcolor{cinzaTabela}\textbf{Bibliografia Complementar:}} \\ \hline
\multicolumn{4}{|p{\dimexpr\linewidth-2\tabcolsep-2\arrayrulewidth\relax}|}{1. PALADINI, Edson Pacheco. Gestão da Qualidade: teoria e prática. Atlas, 2019.} \\ \hline
\end{tabularx}
\end{table}

\vspace{0.3cm}


% ---------------------------------------------------------
% TABELA EMENTA 44: EMPREENDEDORISMO II (PLANO DE NEGÓCIOS)
% ---------------------------------------------------------
\begin{table}[H]
\centering
\small
\begin{tabularx}{\linewidth}{|p{3.4cm}|X|p{3.5cm}|X|}
\hline
\multicolumn{4}{|>{\centering\arraybackslash\cellcolor{ifscgreen!20}}p{\dimexpr\linewidth-2\tabcolsep-2\arrayrulewidth\relax}|}{\textbf{\large Unidade Curricular: EMPREENDEDORISMO II (PLANO DE NEGÓCIOS)}} \\ \hline
\textbf{CH Total*:} & 80 horas & \textbf{Ano / Semestre*:} & 3º Ano (Semestre 1) \\ \hline
\textbf{CH* Prática:} & 40 horas & \textbf{CH EaD*:} & não se aplica \\ \hline
\multicolumn{2}{|l|}{\textbf{CH com Divisão de Turma*:}} & \multicolumn{2}{l|}{não se aplica} \\ \hline
\multicolumn{4}{|p{\dimexpr\linewidth-2\tabcolsep-2\arrayrulewidth\relax}|}{\cellcolor{cinzaTabela}\textbf{Objetivos:}} \\ \hline
\multicolumn{4}{|p{\dimexpr\linewidth-2\tabcolsep-2\arrayrulewidth\relax}|}{
\begin{itemize}[leftmargin=*,noitemsep,topsep=2pt]
    \item Elaborar um Plano de Negócios completo contemplando análise financeira, mercadológica e operacional.
\end{itemize}
} \\ \hline
\multicolumn{4}{|p{\dimexpr\linewidth-2\tabcolsep-2\arrayrulewidth\relax}|}{\cellcolor{cinzaTabela}\textbf{Conteúdos:}} \\ \hline
\multicolumn{4}{|p{\dimexpr\linewidth-2\tabcolsep-2\arrayrulewidth\relax}|}{
\begin{itemize}[leftmargin=*,noitemsep,topsep=2pt]
    \item Estrutura do Plano de Negócios (Sumário Executivo, Análise de Mercado).
    \item Plano Operacional, Financeiro e Estratégico.
    \item Análise de Viabilidade Econômica.
\end{itemize}
} \\ \hline
\multicolumn{4}{|p{\dimexpr\linewidth-2\tabcolsep-2\arrayrulewidth\relax}|}{\cellcolor{cinzaTabela}\textbf{Metodologia de Abordagem:}} \\ \hline
\multicolumn{4}{|p{\dimexpr\linewidth-2\tabcolsep-2\arrayrulewidth\relax}|}{Desenvolvimento orientado de Plano de Negócios em equipe. Banca examinadora.} \\ \hline
\multicolumn{4}{|p{\dimexpr\linewidth-2\tabcolsep-2\arrayrulewidth\relax}|}{\cellcolor{cinzaTabela}\textbf{Bibliografia Básica:}} \\ \hline
\multicolumn{4}{|p{\dimexpr\linewidth-2\tabcolsep-2\arrayrulewidth\relax}|}{1. DORNELAS, José. Empreendedorismo: Transformando Ideias em Negócios. Atlas, 2018.} \\ \hline
\multicolumn{4}{|p{\dimexpr\linewidth-2\tabcolsep-2\arrayrulewidth\relax}|}{\cellcolor{cinzaTabela}\textbf{Bibliografia Complementar:}} \\ \hline
\multicolumn{4}{|p{\dimexpr\linewidth-2\tabcolsep-2\arrayrulewidth\relax}|}{1. BIAGIO, Luiz Antonio. Plano de Negócios: Estratégia para o Sucesso. Manole, 2020.} \\ \hline
\end{tabularx}
\end{table}

\vspace{0.3cm}


% ---------------------------------------------------------
% TABELA EMENTA 45: RESPONSABILIDADE SOCIOAMBIENTAL E SUSTENTABILIDADE
% ---------------------------------------------------------
\begin{table}[H]
\centering
\small
\begin{tabularx}{\linewidth}{|p{3.4cm}|X|p{3.5cm}|X|}
\hline
\multicolumn{4}{|>{\centering\arraybackslash\cellcolor{ifscgreen!20}}p{\dimexpr\linewidth-2\tabcolsep-2\arrayrulewidth\relax}|}{\textbf{\large Unidade Curricular: RESPONSABILIDADE SOCIOAMBIENTAL E SUSTENTABILIDADE}} \\ \hline
\textbf{CH Total*:} & 80 horas & \textbf{Ano / Semestre*:} & 3º Ano (Semestre 1) \\ \hline
\textbf{CH* Prática:} & 30 horas & \textbf{CH EaD*:} & não se aplica \\ \hline
\multicolumn{2}{|l|}{\textbf{CH com Divisão de Turma*:}} & \multicolumn{2}{l|}{não se aplica} \\ \hline
\multicolumn{4}{|p{\dimexpr\linewidth-2\tabcolsep-2\arrayrulewidth\relax}|}{\cellcolor{cinzaTabela}\textbf{Objetivos:}} \\ \hline
\multicolumn{4}{|p{\dimexpr\linewidth-2\tabcolsep-2\arrayrulewidth\relax}|}{
\begin{itemize}[leftmargin=*,noitemsep,topsep=2pt]
    \item Compreender a integração dos critérios ESG na gestão das empresas modernas e sustentabilidade regional.
\end{itemize}
} \\ \hline
\multicolumn{4}{|p{\dimexpr\linewidth-2\tabcolsep-2\arrayrulewidth\relax}|}{\cellcolor{cinzaTabela}\textbf{Conteúdos:}} \\ \hline
\multicolumn{4}{|p{\dimexpr\linewidth-2\tabcolsep-2\arrayrulewidth\relax}|}{
\begin{itemize}[leftmargin=*,noitemsep,topsep=2pt]
    \item Conceitos de Sustentabilidade e Objetivos de Desenvolvimento Sustentável (ODS/ONU).
    \item Práticas ESG (Environmental, Social, and Governance) nas Organizações.
    \item Gestão de Resíduos, Logística Reversa e Economia Circular.
\end{itemize}
} \\ \hline
\multicolumn{4}{|p{\dimexpr\linewidth-2\tabcolsep-2\arrayrulewidth\relax}|}{\cellcolor{cinzaTabela}\textbf{Metodologia de Abordagem:}} \\ \hline
\multicolumn{4}{|p{\dimexpr\linewidth-2\tabcolsep-2\arrayrulewidth\relax}|}{Estudos de casos de empresas sustentáveis e projetos de intervenção. Apresentação final.} \\ \hline
\multicolumn{4}{|p{\dimexpr\linewidth-2\tabcolsep-2\arrayrulewidth\relax}|}{\cellcolor{cinzaTabela}\textbf{Bibliografia Básica:}} \\ \hline
\multicolumn{4}{|p{\dimexpr\linewidth-2\tabcolsep-2\arrayrulewidth\relax}|}{1. KRAEMER, Maria Elisabeth Pereira. Gestão Ambiental e Responsabilidade Social. Atlas, 2012.} \\ \hline
\multicolumn{4}{|p{\dimexpr\linewidth-2\tabcolsep-2\arrayrulewidth\relax}|}{\cellcolor{cinzaTabela}\textbf{Bibliografia Complementar:}} \\ \hline
\multicolumn{4}{|p{\dimexpr\linewidth-2\tabcolsep-2\arrayrulewidth\relax}|}{1. TACHIZAWA, Takeshy. Gestão Ambiental e Responsabilidade Social Corporativa. Atlas, 2019.} \\ \hline
\end{tabularx}
\end{table}

\vspace{0.3cm}


% ---------------------------------------------------------
% TABELA EMENTA 46: OFICINA DE INTEGRAÇÃO III (PROJETO INTEGRADOR)
% ---------------------------------------------------------
\begin{table}[H]
\centering
\small
\begin{tabularx}{\linewidth}{|p{3.4cm}|X|p{3.5cm}|X|}
\hline
\multicolumn{4}{|>{\centering\arraybackslash\cellcolor{ifscgreen!20}}p{\dimexpr\linewidth-2\tabcolsep-2\arrayrulewidth\relax}|}{\textbf{\large Unidade Curricular: OFICINA DE INTEGRAÇÃO III (PROJETO INTEGRADOR)}} \\ \hline
\textbf{CH Total*:} & 40 horas & \textbf{Ano / Semestre*:} & 3º Ano (Semestres 1 e 2) \\ \hline
\textbf{CH* Prática:} & 30 horas & \textbf{CH EaD*:} & não se aplica \\ \hline
\multicolumn{2}{|l|}{\textbf{CH com Divisão de Turma*:}} & \multicolumn{2}{l|}{não se aplica} \\ \hline
\multicolumn{4}{|p{\dimexpr\linewidth-2\tabcolsep-2\arrayrulewidth\relax}|}{\cellcolor{cinzaTabela}\textbf{Objetivos:}} \\ \hline
\multicolumn{4}{|p{\dimexpr\linewidth-2\tabcolsep-2\arrayrulewidth\relax}|}{
\begin{itemize}[leftmargin=*,noitemsep,topsep=2pt]
    \item Integrar saberes científicos e técnicos na solução de problemas reais da região de Garopaba;
    \item Apresentar resultados em formato de artigo científico, relatório técnico e feira acadêmica.
\end{itemize}
} \\ \hline
\multicolumn{4}{|p{\dimexpr\linewidth-2\tabcolsep-2\arrayrulewidth\relax}|}{\cellcolor{cinzaTabela}\textbf{Conteúdos:}} \\ \hline
\multicolumn{4}{|p{\dimexpr\linewidth-2\tabcolsep-2\arrayrulewidth\relax}|}{
\begin{itemize}[leftmargin=*,noitemsep,topsep=2pt]
    \item Metodologia Científica e Diagnóstico de Problemas Locais.
    \item Elaboração do Projeto de Pesquisa/Extensão Aplicada.
    \item Escrita do Relatório Final e Apresentação no Encontro Acadêmico.
\end{itemize}
} \\ \hline
\multicolumn{4}{|p{\dimexpr\linewidth-2\tabcolsep-2\arrayrulewidth\relax}|}{\cellcolor{cinzaTabela}\textbf{Metodologia de Abordagem:}} \\ \hline
\multicolumn{4}{|p{\dimexpr\linewidth-2\tabcolsep-2\arrayrulewidth\relax}|}{Orientação de projetos integradores com encontros presenciais e acompanhamento tutorial. Defesa em banca.} \\ \hline
\multicolumn{4}{|p{\dimexpr\linewidth-2\tabcolsep-2\arrayrulewidth\relax}|}{\cellcolor{cinzaTabela}\textbf{Bibliografia Básica:}} \\ \hline
\multicolumn{4}{|p{\dimexpr\linewidth-2\tabcolsep-2\arrayrulewidth\relax}|}{1. MARCONI, Marina de Andrade; LAKATOS, Eva Maria. Fundamentos de Metodologia Científica. Atlas, 2017.} \\ \hline
\multicolumn{4}{|p{\dimexpr\linewidth-2\tabcolsep-2\arrayrulewidth\relax}|}{\cellcolor{cinzaTabela}\textbf{Bibliografia Complementar:}} \\ \hline
\multicolumn{4}{|p{\dimexpr\linewidth-2\tabcolsep-2\arrayrulewidth\relax}|}{1. SEVERINO, Antônio Joaquim. Metodologia do Trabalho Científico. Cortez, 2017.} \\ \hline
\end{tabularx}
\end{table}

\vspace{0.3cm}


\subsection{Estágio Curricular Supervisionado:}
O Estágio Curricular no Curso Técnico em Administração Integrado ao Ensino Médio é \textbf{não obrigatório}, em consonância com a Resolução CNE/CP nº 1/2021 e o Regulamento Didático-Pedagógico (RDP) do IFSC. O estudante poderá realizar estágio não obrigatório como atividade de prática profissional complementar, devendo haver celebração do Termo de Compromisso de Estágio (TCE) e supervisão docente.

\subsection{Atendimento e Acompanhamento ao Discente:}
O Câmpus Garopaba dispõe de estrutura multiprofissional de apoio aos discentes, composta por:
\begin{itemize}
    \item \textbf{Núcleo de Atendimento às Pessoas com Necessidades Específicas (NAPNE):} Acompanhamento e adaptação curricular para estudantes com deficiência, transtornos globais do desenvolvimento e altas habilidades;
    \item \textbf{Assistência Estudantil (PAE):} Concessão de auxílios financeiros, suporte socioeconômico e acompanhamento de permanência e êxito;
    \item \textbf{Núcleo Pedagógico e de Apoio ao Estudante (NPE):} Atendimento pedagógico, orientação de estudos e mediação escolar.
\end{itemize}

\subsection{Critérios de Aproveitamento de Conhecimentos e Experiências Anteriores:}
O aproveitamento de estudos e a validação de conhecimentos anteriores adquiridos no ensino formal, no trabalho ou em outros espaços de formação obedecerão às disposições do Regulamento Didático-Pedagógico do IFSC (RDP), mediante avaliação teórica/prática por banca examinadora constituída.

\subsection{Certificações Intermediárias:}
Não há previsão de certificação profissional intermediária ao longo do curso. A diplomação técnica integral ocorre mediante a conclusão com aprovação de todas as Unidades Curriculares e cumprimento da carga horária mínima de 3.100 horas-relógio.

% =========================================================
\section{AVALIAÇÃO E ACOMPANHAMENTO DO PPC}
\subsection{Avaliação da Aprendizagem:}
Conforme o Regulamento Didático-Pedagógico (RDP) do IFSC, a avaliação da aprendizagem possui caráter diagnóstico, formativo e processual, pautada em critérios claros de competências desenvolvidas ao longo de cada fase.

\subsection{Autoavaliação do Curso e Atuação da CPA:}
A avaliação do curso é realizada periodicamente em articulação com a Comissão Própria de Avaliação (CPA) do IFSC, comitês pedagógicos e reuniões do Colegiado do Curso, promovendo a melhoria contínua da infraestrutura, matriz curricular e metodologias de ensino.

% =========================================================
\section{INFRAESTRUTURA E RECURSOS}
\subsection{Instalações e Equipamentos:}
O Câmpus Garopaba conta com salas de aula equipadas com recursos multimídia, laboratórios de informática com softwares de gestão, biblioteca com acervo físico e digital, além de espaços para reuniões pedagógicas e atividades esportivas e culturais.

\subsection{Biblioteca e Acervos Digitais:}
Atendimento às normativas do GT de Acervos Digitais e Diretrizes do IFSC, garantindo aos discentes acesso à Plataforma de Livros Digitais (Minha Biblioteca/Pearson), base de dados Target GEDWeb e periódicos científicos.

% =========================================================
\section{CORPO DOCENTE E TÉCNICO-ADMINISTRATIVO}
\subsection{Corpo Docente:}
Composto por professores do quadro permanente do IFSC com qualificação acadêmica de especialização, mestrado e doutorado nas áreas de Administração, Economia, Contabilidade, Direito, Educação e nas disciplinas de Formação Geral (Linguagens, Matemática, Ciências da Natureza e Humanas).

\subsection{Corpo Técnico-Administrativo:}
Suporte por equipe multidisciplinar de Técnicos em Assuntos Educacionais, Assistentes em Administração, Psicólogos, Assistentes Sociais e Técnicos de Laboratório de Informática.

% =========================================================
\section{REFERÊNCIAS NORMATIVAS}
\begin{itemize}
    \item BRASIL. \textit{Lei nº 9.394, de 20 de dezembro de 1996}. Estabelece as diretrizes e bases da educação nacional.
    \item BRASIL. \textit{Lei nº 11.892, de 29 de dezembro de 2008}. Cria a Rede Federal de Educação Profissional, Científica e Tecnológica.
    \item BRASIL. Conselho Nacional de Educação. \textit{Resolução CNE/CP nº 1, de 5 de janeiro de 2021}. Diretrizes Curriculares Nacionais Gerais para a Educação Profissional e Tecnológica.
    \item IFSC. \textit{Resolução CONSUP/IFSC nº 142/2025}. Diretrizes para a organização dos cursos técnicos integrados ao ensino médio no IFSC.
    \item IFSC. \textit{Regulamento Didático-Pedagógico (RDP)}. Aprovado pelo CONSUP/IFSC.
\end{itemize}

\end{document}
